\documentclass[spanish,a4paper,12pt,oneside]{book}
\usepackage{lmodern}
\usepackage{amssymb,amsmath}
\usepackage{ifxetex,ifluatex}
\usepackage{fixltx2e} % provides \textsubscript
\ifnum 0\ifxetex 1\fi\ifluatex 1\fi=0 % if pdftex
  \usepackage[T1]{fontenc}
  \usepackage[utf8]{inputenc}
\else % if luatex or xelatex
  \ifxetex
    \usepackage{mathspec}
  \else
    \usepackage{fontspec}
  \fi
  \defaultfontfeatures{Ligatures=TeX,Scale=MatchLowercase}
\fi
% use upquote if available, for straight quotes in verbatim environments
\IfFileExists{upquote.sty}{\usepackage{upquote}}{}
% use microtype if available
\IfFileExists{microtype.sty}{%
\usepackage{microtype}
\UseMicrotypeSet[protrusion]{basicmath} % disable protrusion for tt fonts
}{}
\usepackage[inner = 3cm, outer = 2.5cm, top = 2.5cm, bottom = 2.5cm]{geometry}
\usepackage{hyperref}
\hypersetup{unicode=true,
            pdftitle={Estudio de los mecanismos físicos asociados con el patrón de onda 3 de la circulación atmosférica del HS},
            pdfauthor={Elio Campitelli},
            pdfborder={0 0 0},
            breaklinks=true}
\urlstyle{same}  % don't use monospace font for urls
\ifnum 0\ifxetex 1\fi\ifluatex 1\fi=0 % if pdftex
  \usepackage[shorthands=off,main=spanish]{babel}
\else
  \usepackage{polyglossia}
  \setmainlanguage[]{spanish}
\fi
\usepackage{natbib}
\bibliographystyle{apalike-es}
\usepackage{longtable,booktabs}
\usepackage{graphicx,grffile}
\makeatletter
\def\maxwidth{\ifdim\Gin@nat@width>\linewidth\linewidth\else\Gin@nat@width\fi}
\def\maxheight{\ifdim\Gin@nat@height>\textheight\textheight\else\Gin@nat@height\fi}
\makeatother
% Scale images if necessary, so that they will not overflow the page
% margins by default, and it is still possible to overwrite the defaults
% using explicit options in \includegraphics[width, height, ...]{}
\setkeys{Gin}{width=\maxwidth,height=\maxheight,keepaspectratio}
\IfFileExists{parskip.sty}{%
\usepackage{parskip}
}{% else
\setlength{\parindent}{0pt}
\setlength{\parskip}{6pt plus 2pt minus 1pt}
}
\setlength{\emergencystretch}{3em}  % prevent overfull lines
\providecommand{\tightlist}{%
  \setlength{\itemsep}{0pt}\setlength{\parskip}{0pt}}
\setcounter{secnumdepth}{5}
% Redefines (sub)paragraphs to behave more like sections
\ifx\paragraph\undefined\else
\let\oldparagraph\paragraph
\renewcommand{\paragraph}[1]{\oldparagraph{#1}\mbox{}}
\fi
\ifx\subparagraph\undefined\else
\let\oldsubparagraph\subparagraph
\renewcommand{\subparagraph}[1]{\oldsubparagraph{#1}\mbox{}}
\fi

%%% Use protect on footnotes to avoid problems with footnotes in titles
\let\rmarkdownfootnote\footnote%
\def\footnote{\protect\rmarkdownfootnote}

%%% Change title format to be more compact
\usepackage{titling}

% Create subtitle command for use in maketitle
\newcommand{\subtitle}[1]{
  \posttitle{
    \begin{center}\large#1\end{center}
    }
}

\setlength{\droptitle}{-2em}
  \title{Estudio de los mecanismos físicos asociados con el patrón de onda 3 de
la circulación atmosférica del HS}
  \pretitle{\vspace{\droptitle}\centering\huge}
  \posttitle{\par}
  \author{Elio Campitelli}
  \preauthor{\centering\large\emph}
  \postauthor{\par}
  \date{}
  \predate{}\postdate{}

\usepackage{setspace}
\setstretch{1.5}
\usepackage{subfig}
\usepackage{hyperref}
\usepackage{marginnote}
\usepackage[nomarkers,figuresonly,nofighead,nolists]{endfloat}
\usepackage{pdflscape}
\DeclareDelayedFloatFlavor{landscape}{figure}
\usepackage[spanish]{todonotes}
\usepackage{wrapfig}
\usepackage{helvet}
\renewcommand{\familydefault}{\sfdefault}
    
\makeatletter
\@ifpackageloaded{endfloat}
{
 \efloat@condopen{fff}
 \pretocmd\@chapter{\immediate\write\efloat@postfff{\noexpand\stepcounter{chapter}}}{}{\fail}
 \AtBeginFigures{\setcounter{chapter}{0}}
}{}
\makeatother
\usepackage{helvet}
\renewcommand{\familydefault}{\sfdefault}

\renewcommand{\figuresection}{Figuras}
\AtBeginDocument{\let\maketitle\relax}
\usepackage[nottoc,notlot,notlof]{tocbibind}
\AtBeginDocument{
\renewcommand{\tablename}{Tabla}
}

\begin{document}
\maketitle

\frontmatter

\begin{titlepage}
    \centering
    \includegraphics[width=0.2\textwidth]{logoUBA}  \hfill \includegraphics[width=0.2\textwidth]{logoDCAO} \par
    \vspace{0.5cm}
    {\scshape\LARGE Universidad de Buenos Aires  \\
    \large Facultad de Ciencias Exactas y Naturales \\
Departamento de Ciencias de la Atmósfera y los Océanos  \par}
    \vspace{0.5cm}
    {\scshape\Large Tesis de Licenciatura en Ciencias de la Atmósfera\par}
    \vspace{1cm}
    {\huge\bfseries Estudio de los mecanismos físicos asociados con el patrón de onda 3 de la circulación atmosférica del HS\par}
    \vspace{1.5cm}
    {\Large Tesista: Elio \textsc{Campitelli} \\
        Directora:  Dra. Carolina Susana \textsc{Vera} \\
        Asistente: Dr. Leandro \textsc{Díaz}
    \par}
    \vfill

% Bottom of the page
    {\large --2018--\par}
\end{titlepage}

\chapter*{Agradecimientos}

A mis directores, Carolina y Leandro, que me guiaron en este
interminable proceso de aprender a investigar que recién comienza.
Espero haber logrado que estén orgullosos de mí y poder seguir
haciéndolo en el futuro.

A los docentes del DCAO, DF y DM que me dieron los conocimientos
necesarios para poder a empezar a aprender a investigar. Ningún libro de
texto tiene el factor humano de esta facultad.

A todas las personas detrás de los paquetes que hicieron posible esta
tesis. Hadley Wickham y los colaboradores de \texttt{ggplot2}, el motor
detrás de todos los gráficos de la tesis y que me permitió
interiorizarme con la gramática de los gráficos. Matt Dowle cuyo
paquete, \texttt{data.table} me permitió analizar millones de datos con
una pequeña y anciana laptop. Y al autor de \texttt{knitr}, Yihui Xie,
quien diseñó una forma elegante de unificar código y texto para una
ciencia reproducible.

A Alexandra Elbakyan. Gracias a ella pude acceder a gran parte del
conocimiento previo que alimentó esta tesis.

A ese elemento emergente que es la comunidad de StackOverflow. Mi
capacidad de resolver los problemas de mi código se reduce
considerablemente cuando no tengo acceso a esa extensión de mi cerebro.

A mi familia. Mi mamá, que me enseñó lo que es el condicionamiento
clásico, mi papá, que un día me sentó en una plaza de San Juan para
hablarme de la gravitación, a Enrique, que me mostró la maleabilidad del
plomo, y a Goyo, que me ayudó a aprender a leer con la ``J'' de
``joder''.

A mis amigos que están ahí hace más de 15 años, acompañando en los
momentos difíciles de la vida y en los momentos fáciles. Son un grupo de
personas increíble que es más que la suma de sus partes.

A Pao, quien hace mi vida infinitamente más rica. No hay nada más
hermoso que aprender y entender (parte) el cosmos con ella. Gracias a
ella soy una mejor persona y un mejor científico.

\newpage

\newgeometry{inner=3.5cm, outer=3.5cm} \chapter*{Resumen}

Este trabajo estudia las características observadas de la onda
cuasiestacionaria 3 (QS3) en el HS a partir del reanálisis del NCEP/NCAR
y analiza preliminarmente experimentos de sensibilidad realizados
utilizando el modelo SPEEDY del ICTP. Los resultados obtenidos muestran
que la QS3 explica una parte importante de la circulación del HS entre
40°S y 60°S y presenta ciclos anuales tanto en su amplitud como en su
fase. Se encontró que la variabilidad estacional de la QS3 se describe
mejor a partir de una clasificación alternativa (EFM, AM, JJ, ASO, ND).
El análisis de sus condiciones de estacionariedad sugiere que la QS3 es
más estacionaria al sur del Índico que al sur del Pacífico. Regresiones
obtenidas a partir de un índice de actividad de la QS3 con diferentes
variables meteorológicas muestran que, según el mes, la QS3 está
asociada tanto con estructuras de ondas planetarias de escala
hemisférica como con trenes de onda más localizados. Esto sugiere un
grado importante de heterogeneidad en los procesos vinculados a su
generación, indicando la posibilidad de que la descripción de la onda 3
a partir de al descomposición de Fourier esté capturando un número de
fenómenos independientes.

Una validación del modelo SPEEDY muestra que éste desarrolla la QS3,
aunque su desempeño fue muy variable según la estación y su inclinación
con la latitud es contraria a la observada en NCEP. El análisis
preliminar de los experimentos de sensibilidad sugiere que las anomalías
zonales de geopotencial en latitudes medias y altas del HS son sensibles
a las condiciones de contorno inferior. En particular la amplitud de la
QS1 disminuye significativamente cuando se eliminan las asimetrías
zonales de la temperatura de la superficie del mar, consistente con
estudios previos que resaltaban el rol de los trópicos en la generación
de este patrón. Asimismo, cuando se elimina la interacción
suelo-atmósfera sobre las zonas continentales, la amplitud de la QS3
aumenta en invierno y disminuye en primavera debido, principalmente, a
cambios en la estacionariedad de la onda. Este resultado permite
concluir que la interacción con el suelo es importante para definir la
localización de la QS3 y no tanto su amplitud, y que las variaciones
estacionales son importantes

\restoregeometry
\setcounter{tocdepth}{3} \tableofcontents

\listoffigures
\newpage

\mainmatter

\chapter{Introducción}\label{introduccion}

La variabilidad del clima en el sur de Sudamérica (SSA) es afectada
tanto por la variabilidad en regiones de latitudes medias y polares
como, remotamente, por las condiciones del clima en regiones tropicales
y extratropicales \citep[ej.][]{Vera2006}. La influencia de éstas
últimas ha sido y es estudiada más extensamente debido a que exhiben
altos niveles de predictibilidad asociados, en escalas estacionales con
el fenómeno del Niño-Oscilación del Sur (ENSO, por sus siglas en inglés)
y en escalas intraestacionales, con la oscilación de Madden-Julian (MJO,
por sus siglas en inglés). Sin embargo, estudios recientes de la
variabilidad climática en las regiones subpolares y polares del
Hemisferio Sur (HS) destacan la actividad del Modo Anular del Sur (SAM,
por sus siglas en inglés) como una fuente de variabilidad climática de
significativa influencia sobre los continentes del HS y en particular en
SSA (ej. \citet{Silvestri2009} en escalas interanuales y
\citet{Alvarez2014} en escalas intraestacionales).

La distribución espacial del SAM en su fase positiva se caracteriza por
una anomalía negativa de presión o altura geopotencial sobre el
continente antártico y anomalías de signo opuesto en latitudes medias.
Esta estructura zonalmente simétrica generalmente está asociada con un
patrón de onda planetario de números de onda zonal entre 3 y 4. La
alternancia del signo entre las anomalías polares y extratropicales del
SAM se asocia con intensificaciones y debilitamiento de los vientos
Oestes que caracterizan las latitudes medias del HS. Si bien hay
evidencia de que los flujos de cantidad de movimiento por las ondas
planetarias de baja frecuencia son importantes en los procesos de
intensificación y debilitamiento de los Oestes
\citep{Lorenz2001, Simpson2013} las causas por las cuales estas
variaciones se asocian con el desarrollo de la QS3 aún no se conocen con
certeza.

Las ondas planetarias de número de onda 1 y 3 (QS1 y QS3) son las
principales asimetrías zonales presentes en el flujo medio del HS
\citep{Loon1972, Trenberth1980a}. Estas ondas zonales tienden a ser
cuasi-estacionarias y exhiben importantes variabilidades temporales en
su amplitud y fase \citep{Loon1972}. En particular, se ha documentado
que la QS3 presenta una estructura barotrópica equivalente con
variabilidad en escalas diarias-semanales \citep{Kidson1988},
estacionales \citep{Mo1985} y más largas \citep{Karoly1989}.
\citet{Mo1985} mostraron que existen ubicaciones preferenciales para los
centros de acción de la QS3. Sin embargo, por la antigüedad de muchos de
estos trabajos, sus climatologías tienen limitaciones inherentes a la
poca cantidad de datos disponible en el HS previo a la era satelital.

\citet{Trenberth1985} mostraron una recurrencia importante en la
ocurrencia de anticiclones de bloqueo simultáneos en diferentes regiones
del HS (sur de Sudamérica, sur de Nueva Zelanda y porción central del
Océano Indico) favorecida por el establecimiento de un patrón de QS3,
aunque dejaron abierta la posibilidad de que éste se trate de un tren de
ondas localizado en vez de una verdadera onda zonal. Desde ese momento
hasta la actualidad, diferentes estudios se concentraron en entender la
influencia de los bloqueos sobre el clima de Sudamérica y su relación
con la QS3 \citep[ej.][]{Rao2004} o de su efecto en el hielo marítimo
\citep[ej.][]{Raphael2007} pero, como se mencionó anteriormente, muy
poca atención ha recibido el estudio de las causas que dan lugar al
establecimiento de este patrón en un primer lugar. \citet{Quintanar1995}
realizaron experimentos de sensibilidad tratando de identificar los
factores importantes en el mantenimiento de la QS1. Encontraron que las
condiciones térmico-orográficas sobre la Antártida no eran suficientes
para explicar la QS1 de latitudes subpolares, por lo que concluyeron que
los forzantes remotos debían jugar un papel importante. \citet{Wang2013}
encontraron que la destrucción y recuperación de la capa de ozono está
asociada a un aumento y disminución de la actividad de onda planetaria
respectivamente, pero su análisis no distingue en la actividad de
distintos números de onda. Por otra parte, \citet{Hobbs2010} pusieron en
duda la utilidad de analizar la circulación del HS en ondas planetarias
zonalmente simétricas y propusieron que las anomalías zonales de
geopotencial están mejor caracterizada por un par de anticiclones al sur
de Nueva Zelanda y al sur de Sudamérica.

\citet{Cai1999} evaluaron la habilidad del modelo CSIRO para representar
la QS3 y encontraron que proporcionaba una representación adecuada de la
misma. \citet{Raphael1998} examinó la QS3 en el modelo CCM del NCAR
(versiones 1 y 3) y encontró diferencias en cómo cada versión simulaba
la QS3, siendo los resultados sensibles al grado de representación del
modelo tanto de la interacción mar-atmósfera como de las condiciones de
hielo marino en las zonas polares. Este es un resultado importante ya
que a pesar de que la QS3 puede desarrollarse solamente por la dinámica
interna de la atmósfera \citep[ej.][]{Simpson2013}, existen evidencias
de que las condiciones superficiales pueden influenciar su actividad.
Asimismo, \citet{Raphael2003} encontró importantes variaciones
interanuales experimentadas por la QS3 entre 1958 y 1996, aparentemente
relacionadas con variaciones en la frecuencia del ENSO, que también
influenciaron las asimetrías del SAM \citep{Fogt2012}.

En consecuencia el objetivo, general de esta tesis es entender los
mecanismos que explican el desarrollo de la QS3 en la circulación del
HS. Los objetivos particulares son

\begin{itemize}
\tightlist
\item
  caracterizar la climatología de la QS3;
\item
  explorar la influencia de las condiciones oceánicas superficiales
  tanto en los trópicos como en los extratrópicos sobre la actividad de
  la QS3;
\item
  explorar la sensibilidad de la QS3 en general a las condiciones
  superficie.
\end{itemize}

\chapter{Datos y Metodologías}\label{datos-y-metodologias}

En este capítulo se describen los datos utilizados y las metodologías
aplicadas en la presente investigación.

\section{Datos y modelo}\label{datos-y-modelo}

Se analizaron datos mensuales de altura geopotencial, temperatura,
viento zonal, viento meridional y función corriente provenientes del
Reanálisis NCEP/NCAR \citep{Kalnay1996} (de aquí en adelante datos NCEP)
entre enero de 1985 y diciembre de 2015. Los mismos poseen una
resolución de 2,5° en longitud y latitud y 17 niveles verticales entre
1000hPa y 10hPa. A partir de los datos mensuales se calcularon las
medias estacionales definidas a partir de las estaciones climatológicas
del HS.

Se utilizó el modelo SPEEDY del ICTP para realizar corridas de
sensibilidad de la QS3 a diferentes condiciones (descriptas en la
\autoref{experimentos-de-sensibilidad}). El modelo SPEEDY
\citep{Molteni2003, Kucharski2006} es un modelo de complejidad
intermedia basado en ecuaciones primitivas espectrales y
parametrizaciones simplificadas. En su versión 41, posee una resolución
horizontal espectral de T30 (3,75° en longitud y latitud) y 8 niveles
verticales entre 925hPa y 30hPa. El modelo incluye parametrizaciones de
las diferentes formas de convección, radiación, flujos y difusión
vertical.

Una primera gran limitación de SPEEDY es su pobre representación de la
estratósfera. El más alto de sus 8 niveles está en 30hPa y al ser la
tapa del modelo, que tiene una ``esponja'' para evitar la propagación de
ondas de gravedad, no es un nivel con información confiable. Esto limita
seriamente la posibilidad de describir numéricamente lo que ocurre en la
estratosfera como, por ejemplo, la dinámica del vórtice polar.

Sin embargo por su bajo costo computacional y su buen desempeño para
simular ciertas características globales de interés para Sudamérica
\citep{Barreiro2014} se decidió utilizarlo en este trabajo.

\section{Metodología}\label{metodologia}

\section*{Perturbaciones}

Muchas variables atmosféricas varían con la longitud en menor medida que
con la latitud o la altura. Resulta entonces natural descomponer
cualquier variable \(\phi_{(x, y, z, t)}\) en un promedio zonal y las
desviaciones con respecto al mismo según:

\[
\phi_{(x, y, z, t)} = [\phi]_{(y, z, t)} + \phi_{(x, y, z, t)}^*
\]

donde los corchetes indican el promedio zonal y el asterisco indica el
desvío con respecto al mismo. \([\phi]\) representa la condición
simétrica zonal de la variable en cuestión y es independiente de la
longitud mientras que \(\phi^*\) representa su componente asimétrica
(que será estacionaria en el caso de considerar el valor medio
climatológico).

Se denomina ondas cuasiestacionarias (QS) a la componente asimétrica de
la circulación media climatológica. Las mismas generan flujos
meridionales de calor y cantidad de movimiento, principalmente en el HN
en invierno \citep{James}, y en gran medida son forzadas por forzantes
superficiales como orografía y contrastes de temperatura así como por
interacciones con los transcientes \citep{Rao2004}.

En este trabajo las anomalías zonales se calcularon para cada círculo de
latitud como la diferencia entre el valor de la variable en una
determinada longitud y su promedio zonal. También se calcularon campos
de anomalía zonal de geopotencial eliminando la influencia de las ondas
planetarias 1 y 2. Para esto se reconstruyeron los campos de anomalía
zonal de geopotencial de QS1 y QS2 mediante Fourier y luego se restaron
al campo de anomalía zonal total.

\section*{Fourier}

La descomposición de Fourier permite representar una serie de datos como
una sumatoria de ondas sinusoidales de distinto número de onda, amplitud
y fase. Adicionalmente, es posible calcular la proporción de varianza
total explicada por cada onda individual~(\(r^2\)) \citep{Wilks2011}.
Para describir las ondas planetarias se aplicó esta metodología a cada
círculo de latitud del campo de geopotencial usando la función
\texttt{FitQsWave()} del paquete \texttt{metR} \citep{R-metR} de manera
de obtener un valor de la amplitud, la fase y el \(r^2\) para los
números de onda 1 a 4. Dado que la fase es una variable circular, su
estudio requiere un tratamiento especial. Para el cálculo de
estadísticas circulares se se utilizó el paquete \texttt{circular}
\citep{R-circular}.

La descripción de las QS a partir de Fourier puede realizarse ya sea
descomponiendo los campos medios climatológicos (de aquí en adelante
\emph{método AM}, por Amplitud de la Media) o calculando el promedio de
la descomposición de los campos instantáneos --mensuales, en este caso--
(de aquí en adelante \emph{método MA}, por Media de la Amplitud
instantánea). Resulta evidente que a partir del método AM se obtiene una
descripción de las ondas presentes en los campos medios y que sobreviven
a la interferencia destructiva del promedio temporal, brindando
información sobre las ondas estacionarias. El método MA, en cambio,
indica las propiedades medias de las ondas planetarias
independientemente de su fase, mezclando información de las anomalías
zonales instantáneas que no necesariamente son estacionarias.

Utilizando los datos NCEP, no se encontraron diferencias importantes
entre ambas metodologías, por lo que se muestran en general los
resultados sólo del método AM (excepto que se indique lo contrario).

\section*{Wavelets}

El análisis de wavelets permite también describir oscilaciones
localizadas en el dominio \citep{Torrence1998}. En el contexto de las
QS, en vez de obtener una amplitud y fase por cada círculo de latitud
como se obtiene con el método de Fourier, con el método de wavelets se
obtiene una amplitud local que depende de la longitud. En la ciencias
atmosféricas el método de wavelets se usa extensivamente en el análisis
de periodicidades temporales \citep[ej.][]{Raphael2004, Kinnard2011} y
de procesamiento de imágenes \citep[ej.][]{Desrochers1999}, pero existen
pocos estudios que lo apliquen a dominios espaciales
\citep[ej.][]{Pinault2016} aunque sí hay aplicaciones de este tipo en la
literatura, por ejemplo, de ecología \citep[ej,][]{Mi2005}.

Para el cálculo de wavelets se utilizó la función
\texttt{WaveletTransform()} del paquete \texttt{WaveletComp}
\citep{R-WaveletComp} que utiliza un wavelet de Morlet.

\section*{Índice de QS3}

En esta tesis se elaboró un índice de actividad de la QS3 a partir de la
amplitud de Fourier promediada entre 65°S y 40°S de latitud y entre
700hPa y 100hPa. Existen múltiples índices para describir la QS3 en la
literatura. Algunos son definidos a partir de las anomalías zonales de
geopotencial en puntos fijos cercanos a donde climatológicamente se dan
los máximos de geopotencial asociados a la QS3
\citep[ej.][\citet{Cai1999} y \citet{Raphael2004}]{Mo1985}. Aunque cada
trabajo utilizó puntos ligeramente distintos a causa de las distintas
climatologías analizadas. Además, \citet{Raphael2004} se basó en
promedios trimestrales en vez de medias mensuales. Por otra parte,
\citet{Yuan2008}, generó una serie temporal de actividad de la QS3 a
partir de la primer componente principal del campo de viento meridional
superficial, el cual resulta en un patrón de onda 3 consistente con
climatologías previas.

La principal limitación de las metodologías basadas en patrones
estacionarios, como puntos fijos o componentes principales, es que no
permiten cambios de fase en la onda planetaria. En particular, no
capturan correctamente el corrimiento estacional en la posición de la
QS3 \citep{Loon1972} ni permiten capturar casos de actividad de onda
intensa pero alejada de las zonas climatológicamente activas.
\citet{Irving2015} construyeron un índice de actividad de onda
planetaria a partir de la transformada de Hilbert que no tiene estas
limitaciones pero no distingue entre la actividad de distintos números
de onda.

\section*{Interferencia constructiva y destructiva}

En la \autoref{fase} se estimó, para cada mes, la proporción de casos de
interferencia destructiva a partir del cómputo de la frecuencia de casos
donde la diferencia de fase con respecto a la fase media es de entre 40°
y 80°. Dicho criterio teórico surge de considerar la suma de dos ondas
sinusoidales de igual amplitud

\[
\cos\left (k\phi \right) + \cos(k(\phi - \alpha)) = 2\cos\left( \frac{k\alpha}{2} \right)\cos\left( k\phi - \frac{\alpha}{2}\right) 
\] donde \(k\) es el número de onda, \(\phi\) la latitud, y \(\alpha\)
es la diferencia de fase entre las ondas. El primer término
multiplicativo del lado derecho es la amplitud de la nueva onda, la cual
depende de la diferencia de fase. Si la misma es menor a 1/3 o mayor a
2/3 de longitud de onda (40° y 80° respectivamente en el caso de la onda
3), el valor absoluto de la amplitud es mayor que 1 y la sumatoria de
ondas es constructiva. De lo contrario, la onda obtenida tiene una
amplitud menor que las originales y la interferencia es destructiva.

En la \autoref{estaciones} se utilizó análisis de componentes
principales para categorizar cada mes del año según las características
típicas de la QS3. Se realizó mediante la función \texttt{EOF()} del
paquete \texttt{metR} \citep{R-metR}, la cual descompone una matriz de
datos en valores singulares. Este método es preferible al cálculo de
autovalores y autovectores sobre la matriz de correlación o covarianza
ya que es numéricamente más estable. Los datos de geopotencial fueron
pesados por la raíz cuadrada del coseno de la latitud.

\section*{Herramientas dinámicas}

La influencia tropical en la variabilidad del clima extratropical se
exploró a través del análisis de las condiciones de propagación
meridional de ondas de Rossby barotrópicas. La teoría lineal establece
que esta propagación depende de las condiciones del entorno
\citep{James}. En particular, el número de onda meridional está dado por

\[
l = \pm \sqrt{\frac{\eta_{y}}{U} - k^2}
\]

donde \(\eta_{y}\) es el gradiente meridional de vorticidad absoluta,
\(U\) es la velocidad zonal, ambos del estado básico y \(k\) es el
número de onda zonal. La propagación meridional sólo es posible si \(l\)
es real, lo que requiere que

\[
\frac{\eta_{y}}{U} = K_s \ge k 
\] donde \(K_s\) es el número de onda estacionario, que debe ser mayor
que el número de onda zonal para permitir la propagación meridional.

El flujo de actividad de onda es un vector que, bajo ciertas
suposiciones, es paralelo a la velocidad de grupo de las ondas de
Rossby, lo cual permite cuantificar la dispersión meridional y zonal de
las mismas \citep{James}. En este estudio, estos flujos se calcularon a
partir de las anomalías zonales de función corriente como

\[
\begin{aligned}
F_\lambda &= \frac{p}{2000a^2\cos\phi}\left[ \left( \frac{\partial \psi^*}{\partial \lambda} \right)^2 - \psi^*\frac{\partial^2 \psi^*}{\partial \lambda^2}  \right] \\
F_\phi &= \frac{p}{2000a^2} \left( \frac{\partial \psi^*}{\partial \lambda}\frac{\partial \psi^*}{\partial \phi}  - \psi^* \frac{\partial^2 \psi^*}{\partial \lambda \partial \phi} \right) 
\end{aligned}
\]

donde \(p\) es la presión, \(\phi\) la latitud, \(\lambda\) la longitud,
\(a\) es el radio de la Tierra (tomado como 6371km) y \(\psi^*\) es la
anomalía zonal de la función corriente \citep{Vera2004}.

\chapter{Climatología observada}\label{climatologia-observada}

En este capítulo se presentan campos medios y anomalías zonales de
altura geopotencial, temperatura, viento zonal, viento meridional,
función corriente, número de onda estacionario y gradiente meridional de
vorticidad absoluta, como introducción general al estado medio de la
atmósfera sobre el cual se desarrollan las ondas estacionarias. Luego se
analizan los campos de amplitud y varianza explicada por las distintas
QS.

\section{Altura geopotencial}\label{altura-geopotencial}

\begin{landscape}\begin{figure}

{\centering \includegraphics[width=247mm,]{fig/tesis/gh-ncep-1} 

}

\caption{Z (mgp). Contornos cada 250 mgp (NCEP).}\label{fig:gh-ncep}
\end{figure}
\end{landscape}

El campo de altura geopotencial media (Z, \autoref{fig:gh-ncep}) muestra
una estructura marcadamente zonal en todos los niveles y estaciones. En
verano el gradiente meridional de Z es máximo en 200hPa, reduciéndose en
500hPa y por encima de 100hPa. En 50hPa el gradiente es prácticamente
nulo y en niveles superiores, éste se invierte en comparación a los
inferiores (no se muestra). En otoño el máximo de gradiente todavía se
da en 200hPa, pero continúa siendo intenso en niveles superiores. En
invierno y primavera, el mayor gradiente se da en 50hPa y es mucho más
intenso que los observados en los demás niveles o estaciones. En
contraste con el resto de los niveles, 50hPa y 100hPa tienen mucha más
variabilidad estacional.

El aumento del gradiente meridional de geopotencial en invierno y
primavera en niveles altos está relacionado con la generación del
vórtice polar que aísla las latitudes polares de las latitudes medias.
En 200hPa, en cambio, es evidente el gradiente asociado con el jet
subtropical, que es más intenso en invierno y más débil en verano.

En la \autoref{fig:sd-gh-ncep} se muestra, para cada latitud y mes, el
desvío estándar de Z con respecto a la media zonal (\(\sigma_z\)). En
todos los niveles la variación meridional es muy similar, con un máximo
principal al rededor de 60°S durante todo el año y un máximo secundario
en 30°S que aparece sólo entre junio y septiembre. El ciclo anual de
\(\sigma_z\) también es similar entre niveles, aunque con ligeros
corrimientos de fase. En 100hPa el máximo absoluto se da en octubre,
mientras que en los niveles más bajos, éste se da en agosto.

\begin{figure*}
\includegraphics[width=155mm, ]{fig/tesis/sd-gh-ncep-1} \caption{Desvío estándar de Z (mgp) por círculo de latitud (NCEP).}\label{fig:sd-gh-ncep}
\end{figure*}

\begin{landscape}\begin{figure}

{\centering \includegraphics[width=247mm,]{fig/tesis/ghz-ncep-1} 

}

\caption{Z* (mgp) (NCEP).}\label{fig:ghz-ncep}
\end{figure}
\end{landscape}

Las anomalías zonales de geopotencial (Z*, \autoref{fig:ghz-ncep})
muestran una preponderancia de la onda 1 (QS1) con una amplitud máxima
en la estratósfera de primavera. Pueden diferenciarse dos QS1 distintas;
una centrada en \textasciitilde{}60°S y con el centro anticiclónico
alrededor de la línea de fecha, y la otra centrada en 75°S sobre la
costa del continente antártico y el centro anticiclónico entre 120 y
60°O. \citet{Quintanar1995} concluyeron que la primer QS1 está asociada
principalmente a forzantes de latitudes bajas mientras que la segunda
responde a la orografía del continente antártico.

En latitudes tropicales, en verano hay una anomalía negativa sobre el
Pacífico este con máxima amplitud en 200hPa que está presente en las
otras estaciones con menor intensidad. Sobre Sudamérica, en verano y
primavera en ese mismo nivel aparece un centro anticiclónico (alta de
Bolivia) con un centro ciclónico al noroeste y otro al noreste (Baja del
Nordeste). Estas anomalías son características del Sistema Monzónico
Sudamericano \citep{Vera2006}.

En la \autoref{fig:ghz-ncep-corte60} se muestra un corte zonal en 60°S
de Z*. Se aprecia la coherencia vertical de la QS1 y es evidente la
inclinación hacia el oeste con la altura en todas las estaciones salvo
en verano, en coincidencia con lo encontrado por \citet{Quintanar1995a}.
\citet{Karoly1985} describió una estructura barotrópica equivalente
tanto en verano como en invierno en 55°S, pero esto se debe a la falta
de niveles verticales por encima de los 100hPa disponibles en ese
momento que es donde la inclinación es más importante.

La inclinación hacia el oeste de las perturbaciones de geopotencial con
la latitud (\autoref{fig:ghz-ncep}) y con la altura
(\autoref{fig:ghz-ncep-corte60}) indica que las mismas están asociadas
con transporte perturbado hacia el polo tanto de cantidad de movimiento
zonal como de temperatura \citep{James}. En verano las anomalías zonales
tienen una estructura barotrópica equivalente y carecen de inclinación
en la horizontal.

\begin{figure*}
\includegraphics[width=155mm, ]{fig/tesis/ghz-ncep-corte60-1} \caption{Corte zonal de Z* (mgp) en 60°S (NCEP).}\label{fig:ghz-ncep-corte60}
\end{figure*}

\section{Temperatura}\label{temperatura}

\begin{landscape}\begin{figure}

{\centering \includegraphics[width=247mm,]{fig/tesis/t-ncep-1} 

}

\caption{Temperatura media (°C). Contornos cada 5°C (NCEP).}\label{fig:t-ncep}
\end{figure}
\end{landscape}

La distribución horizontal de la temperatura media
(\autoref{fig:t-ncep}), al igual que la altura geopotencial media, tiene
una estructura principalmente zonal en todos los niveles y estaciones.
Por debajo de los 200hPa, donde el gradiente meridional de temperatura
es mínimo, la temperatura disminuye con la latitud en todas las
estaciones. Por encima de este nivel, en cambio, en verano la
temperatura crece con la latitud, y en el resto de las estaciones
muestra un máximo centrado en 60°S en otoño y en 45°S en invierno y
primavera. En el nivel de 850hPa se ven las asimetrías zonales más
importantes asociadas con los contrastes de temperatura entre continente
y océanos, como el máximo sobre Australia en verano.

En 300hPa, el gradiente meridional de temperatura en latitudes medias
tiene un importante ciclo anual con máximo en invierno y mínimo en
verano. En niveles inferiores, el ciclo anual es menos marcado y más
dependiente de la latitud. En 500hPa, en 150°E, 30°S se da un máximo del
gradiente meridional de temperatura en invierno y un mínimo en verano,
mientras que en 50°S el ciclo es el inverso.

\begin{figure*}
\includegraphics[width=155mm, ]{fig/tesis/t-ncep-corte-1} \caption{Media zonal de la temperatura (°C) para cada nivel vertical y latitud (NCEP).}\label{fig:t-ncep-corte}
\end{figure*}

El promedio zonal de la temperatura se muestra en la
\autoref{fig:t-ncep-corte}. Al norte de los 45° la temperatura decrece
con la altura por debajo de los 100hPa aproximadamente, donde alcanza un
mínimo que marca la tropopausa, y crece por encima durante todo el año.
La altura del mínimo de temperatura varía mucho estacionalmente al sur
de los 45°, siendo mínima en verano (300hPa) y máxima en invierno y
otoño (30hPa). Hay que tener en cuenta, sin embargo, que la tropopausa
no está bien definida en el invierno antártico
\citep{Court1942, Zangl2001} por lo que el uso de la tropopausa térmica
no es conveniente. \citet{Zangl2001}, utilizando datos de ERA para el
período 1979-93, encontraron que la tropopausa térmica varía entre los
320hPa en enero y los 170hPa en agosto, aunque advierten que el uso del
criterio térmico es ``problemático'' en el invierno antártico.

\begin{landscape}\begin{figure}

{\centering \includegraphics[width=247mm,]{fig/tesis/tz-ncep-1} 

}

\caption{Anomalía zonal de temperatura (°C) (NCEP).}\label{fig:tz-ncep}
\end{figure}
\end{landscape}

En la \autoref{fig:tz-ncep} se muestran las anomalías zonales de
temperatura media. En 850hPa se aprecia el efecto del contraste de
temperatura entre el suelo y el mar. Se observan anomalías positivas
sobre los continentes y negativas sobre los océanos en todas las
estaciones, aunque más intensas en verano y primavera. En niveles más
altos éstas pierden su intensidad pero reaparecen en 100hPa con signo
invertido. Estas características tienen su correlato en la altura
geopotencial (\autoref{fig:ghz-ncep}) y corresponden a circulaciones de
tipo monzónico.

En invierno y primavera, los niveles altos están dominados por una QS1
con máximo en el sur de Australia y mínimo en el Atlántico sur. En
niveles más bajos, esta onda disminuye su amplitud y se defasa hacia el
este y queda casi en cuadratura con la onda de niveles altos presentando
un máximo en 850hPa en Antártida occidental, como se observa en el corte
zonal de la anomalía zonal de temperatura en 60°S
(\autoref{fig:t-ncep-corte60}). En otoño la QS1 también está presente,
pero con amplitud muy reducida y maximizando en niveles medios. En
verano ésta desaparece por encima de 100hPa.

\begin{figure*}
\includegraphics[width=155mm, ]{fig/tesis/t-ncep-corte60-1} \caption{Corte zonal de anomalía zonal de temperatura (°C) en 60°S (NCEP).}\label{fig:t-ncep-corte60}
\end{figure*}

Las anomalías por debajo de 300hPa mantienen su intensidad durante todo
el año, aunque tienen más extensión en invierno y primavera en
comparación con verano y otoño y son barotrópicas equivalentes. Por
encima de 300hPa, en cambio, se observa un importante ciclo anual con
máximo en primavera y mínimo en verano con importante inclinación hacia
el oeste con la altura.

\section{Viento zonal}\label{viento-zonal}

\begin{figure*}
\includegraphics[width=155mm, ]{fig/tesis/u-ncep-corte-1} \caption{Media zonal del viento zonal (m/s) para cada nivel vertical y latitud (NCEP).}\label{fig:u-ncep-corte}
\end{figure*}

La media zonal del viento zonal (\autoref{fig:u-ncep-corte}) muestra dos
máximos, uno en latitudes medias en 200hPa y otro en latitudes polares
en la estratósfera, correspondientes al jet subtropical y subpolar,
respectivamente. El primero está presente durante todo el año, aunque
con mayor intensidad y corrido hacia latitudes más ecuatoriales en
invierno y primavera. El segundo está presente principalmente en
invierno y primavera, e incipiente en otoño. Además, en la estratósfera
se observan vientos del este en latitudes bajas que son más intensos en
verano y otoño.

\begin{landscape}\begin{figure}

{\centering \includegraphics[width=247mm,]{fig/tesis/u-ncep-1} 

}

\caption{Viento zonal medio (m/s) (NCEP).}\label{fig:u-ncep}
\end{figure}
\end{landscape}

En la \autoref{fig:u-ncep} se muestra el viento zonal medio. Se observa,
principalmente en invierno y primavera por encima de 100hPa, que el jet
subpolar es más intenso al sur de África, donde además se encuentra en
una latitud más ecuatorial que en la región del Pacífico. El jet
subtropical también tiene un máximo al sur de África y otro al norte de
Nueva Zelanda --especialmente en invierno--, donde además se produce una
bifurcación del jet. Se trata de una región de persistentes y frecuentes
bloqueos \citep[ej.][]{Trenberth1985}.

Esta bifurcación del jet sobre Nueva Zelanda se evidencia en el campo de
anomalías zonales de viento zonal (\autoref{fig:uz-ncep}) como una
anomalía negativa sobre la isla acompañada por anomalías positivas al
norte y sur. Este campo también presenta varios pares de QS1
antisimétricos respecto a 60°S. Estas anomalías se corresponden con la
variación meridional del jet observada en la \autoref{fig:u-ncep} y son
consistentes con la QS1 de geopotencial observada en la
\autoref{fig:ghz-ncep}. Por el balance de viento geostrófico, centros
anticiclónicos de Z* están flanqueados por anomalías zonales negativas
de viento zonal al norte y positivas al sur y viceversa para los centros
ciclónicos. Además, como el viento zonal no depende del parámetro de
Coriolis como Z, esto permite la identificación de QS en regiones
tropicales. Un ejemplo de esto son la alternancia de centros de
diferente signo en las zonas tropicales en verano, que describen las
ondas de Rossby estacionarias asociadas con los sistemas monzónicos de
Sudamérica y el norte de Australia. En particular, se distingue entre
300hPa y 100hPa sobre el Pacífico ecuatorial una zona de anomalías del
viento zonal positivas al este y negativas al oeste, que implica
divergencias en niveles altos (compensada con convergencias en niveles
bajos que no se muestra).

\begin{landscape}\begin{figure}

{\centering \includegraphics[width=247mm,]{fig/tesis/uz-ncep-1} 

}

\caption{Anomalía zonal de viento zonal (m/s) (NCEP).}\label{fig:uz-ncep}
\end{figure}
\end{landscape}

\section{Viento meridional}\label{viento-meridional}

\begin{figure*}
\includegraphics[width=155mm, ]{fig/tesis/v-ncep-corte-1} \caption{Media zonal del viento meridional (m/s) (NCEP).}\label{fig:v-ncep-corte}
\end{figure*}

El corte vertical-meridional del promedio zonal viento meridional medio
(V, \autoref{fig:v-ncep-corte}) muestra los máximos tropicales presentes
en superficie y altura en todas las estaciones, relacionados con la
circulación de Hadley. En verano, la rama ascendente se encuentra en el
HS y se tiene convergencias en niveles bajos y divergencias en niveles
altos. En invierno, en cambio, sólo se ve la rama descendente, mucho más
intensa que en verano, que genera convergencias en niveles altos y
divergencias en niveles bajos al rededor de los 30°S.

Presente durante todo el año, también se observa un máximo de vientos
del sur en la costa antártica. Los mismos son evidencia de los vientos
catabáticos antárticos producidos por una capa muy estable cerca de
superficie y la consistente inclinación de la topografía del continente
\citep{King1997}. Sin embargo, los datos allí pueden tener limitaciones
por la falta de observaciones y la pobre representación de la orografía
en los modelos.

\begin{landscape}\begin{figure}

{\centering \includegraphics[width=247mm,]{fig/tesis/v-ncep-1} 

}

\caption{Viento meridional medio (m/s) (NCEP).}\label{fig:v-ncep}
\end{figure}
\end{landscape}

Dado que la media zonal de V es nula en gran parte del dominio, el campo
de anomalías zonales de V es virtualmente idéntico al campo de V medio.
Por lo tanto, sólo se muestra el campo horizontal de V en la
\autoref{fig:v-ncep}. Consistente con los campos de Z*
(\autoref{fig:ghz-ncep}), en niveles altos se observa una QS1 que
alcanza su máximo en la estratósfera de primavera. En invierno entre
500hPa y 100hPa, existe evidencia de un tren de ondas de Rossby que se
propaga desde el Índico occidental sudeste llegando a su máxima latitud
en 150°O donde comienza a propagarse hacia el norte hasta llegar al sur
de Sudamérica. Este tren de ondas puede identificarse en el campo de Z*
(\autoref{fig:ghz-ncep}), pero con mayor dificultad debido a la
presencia de la QS1 y a la dependencia de Z* con el parámetro de
Coriolis. El tren de ondas en V se distingue en la troposfera alta,
también en primavera y con menor intensidad en verano y otoño.

Las anomalías zonales del viento meridional también permiten distinguir
otras características cuasiestacionarias del clima, como aquellas
asociadas con los monzones. En el invierno, en los trópicos se puede
observar la anomalía positiva en 850hPa en la costa oeste de África
asociada con el flujo hacia el monzón de la India. En altura, el monzón
de la India se evidencia en esa estación como anomalía de viento hacia
el sur producto de la divergencia de niveles altos generada por la
convección anómala. Por otra parte, en verano se evidencia en la
troposfera alta sobre Sudamérica tropical anomalías positivas sobre el
centro-este del continente y negativas en los océanos subyacentes. Tal
patrón de ondas estacionario se relaciona con la presencia del Alta de
Bolivia y las vaguadas estacionarias a sus lados, siendo la del este
típicamente llamada la Baja del Nordeste \citep{Vera2006}.

\section{Función corriente}\label{funcion-corriente}

\begin{landscape}\begin{figure}

{\centering \includegraphics[width=247mm,]{fig/tesis/psi-ncep-1} 

}

\caption{Función corriente media en $\sigma = 0,2101$ (contornos cada $2\times10^{-11}m^2/s$), anomalía zonal de función corriente (sombreado,  $1\times10^{-9}m^2/s$) y flujos de actividad de onda medios (NCEP).}\label{fig:psi-ncep}
\end{figure}
\end{landscape}

Para analizar la influencia tropical en la circulación no es posible
usar la altura geopotencial, ya que el balance geostrófico pierde
validez cerca del ecuador. Por lo tanto, es útil analizar el campo de
función corriente (\(\psi\)). Los datos NCEP proveen esta variable en
niveles sigma (\(\sigma\)) en vez de presión. En la
\autoref{fig:psi-ncep} se muestra \(\psi\) en el nivel
\(\sigma = 0,2101\), que equivale aproximadamente a 250hPa. Además del
campo medio, se muestran las anomalías zonales de esta variable en
sombreado y los flujos de actividad de onda en flechas.

Las características de los campos de \(\psi\), tanto el total como las
anomalías zonales, son similares a los de Z y Z* respectivamente, con
una estructura eminentemente zonal y un aumento del gradiente meridional
en invierno y primavera y los mismos centros de anomalías. La principal
diferencia es, además del cambio de signo dada por la dependencia de Z
con el parámetro de Coriolis, que los patrones presentes en las
latitudes tropicales se ven con mayor magnitud que los de latitudes
medias y altas.

Los flujos de actividad de onda en verano muestran transporte de energía
desde el Pacífico este hacia el sur de África pasando por el Atlántico
que se sostiene durante todo el año con menor intensidad. Desde ese
lugar también se observa transporte de energía hacia el hemisferio
norte, que se junta con otra región de flujos intensos que viene desde
el Pacífico oeste. Sobre el Índico, los flujos son de mayor magnitud en
invierno, transportando energía hacia el sur.

\section{Propagación Meridional de Ondas de
Rossby}\label{propagacion-meridional-de-ondas-de-rossby}

\begin{landscape}\begin{figure}

{\centering \includegraphics[width=247mm,]{fig/tesis/etady-ncep-1} 

}

\caption{Gradiente meridional de vorticidad absoluta ($1\times10^{11}(ms)^{-1}$). La línea negra marca regiones con valores negativos.}\label{fig:etady-ncep}
\end{figure}
\end{landscape}

La \autoref{fig:etady-ncep} muestra el campo de \(\eta_y\). Resalta una
región de valores negativos centrada sobre Nueva Zelanda, en invierno
entre 300hPa y 200hPa, coincidiendo con la región de bloqueos. Está
flanqueada por el jet subtropical intenso y el jet subpolar más al sur,
dando lugar a gradientes meridionales de viento zonal negativos más
intensos que \(\beta\). Otras regiones con valores negativos se
distinguen en la costa antártica.

Los valores de \(\eta_{y}\) negativos impiden la propagación meridional
de ondas de Rossby barotrópicas. Esta figura reproduce y extiende el
resultado de \citet{Berbery1992} (su Figura~3) que utilizaron 5 años de
análisis objetivo del Centro Europeo de Predicción a Plazo Medio
(ECMWF). Asimismo, como se describió en la \autoref{metodologia}, aún
con \(\eta_{y}\) positivo, las ondas de Rossby barotrópicas sólo se
pueden propagar si su número de onda zonal es menor que el número de
onda estacionario \citep{James}. En la \autoref{fig:ks-ncep} se muestra
el número de onda estacionario para el nivel de 200hPa. Entre otoño y
primavera se destaca claramente la región del Pacífico Oeste en la cual
la propagación meridional está inhibida. Además, en todas las
estaciones, las ondas cortas no pueden propagarse meridionalmente en
latitudes altas. La \autoref{fig:ks-ncep} muestra que las ondas de
Rossby barotrópicas con k menores o iguales a 3 pueden propagarse
meridionalmente en las cuatro estaciones hasta aproximadamente los 60°S,
excepto en la franja de longitudes entre 60°E y 120°O, donde quedan
atrapadas al sur de 45°S (excepto en verano).

\begin{landscape}\begin{figure}

{\centering \includegraphics[width=247mm,]{fig/tesis/ks-ncep-1} 

}

\caption{Número de onda estacionario en 200hPa (NCEP).}\label{fig:ks-ncep}
\end{figure}
\end{landscape}

\section{Ondas Quasiestacionarias}\label{ondas-quasiestacionarias}

La \autoref{fig:r2-ncep} muestra el corte vertical-meridional de \(r^2\)
de Fourier (definido en la \autoref{metodologia}) para las ondas
estacionarias 1 a 4 de manera de analizar la importancia relativa de
cada una con respecto a la variabilidad total. Consistente con lo
encontrado previamente en el campo de Z*, la QS1 explica la mayor parte
de la variabilidad en todo el dominio al sur de los 45°S. La QS2 es
preponderante en la estratósfera ecuatorial, en la costa antártica y
alrededor de 35°S, donde es el modo dominante en toda la columna de aire
en verano. La QS3, a diferencia de las ondas anteriores, es importante
en una región reducida. Explica una parte substancial de la varianza en
niveles bajos al rededor de los 45°S y mayormente en invierno. La QS4
explica muy poca varianza a excepción de cerca de superficie entre 15°S
y 30°S. Ondas más cortas son aún menos importantes (no se muestra).

\begin{landscape}\begin{figure}

{\centering \includegraphics[width=247mm,]{fig/tesis/r2-ncep-1} 

}

\caption{$r^2$ de Fourier para números de onda 1 a 4 (NCEP).}\label{fig:r2-ncep}
\end{figure}
\end{landscape}

El \(r^2\) permite analizar la importancia relativa de cada modo con
respecto a la variabilidad total pero la amplitud, cuyo corte
vertica-meridional se muestra en la \autoref{fig:ampl-ncep}) permite
estimar la magnitud de cada onda independientemente de las demás. Las
diferencias entre los campos de \(r^2\) y los de amplitud son evidentes
comparando esta figura con \autoref{fig:r2-ncep} (notar la escala
logarítmica en los colores). La amplitud de la QS1 muestra un máximo
bien definido centrado en 60°S que en verano se encuentra en niveles más
bajos que en las otras estaciones. También exhibe un máximo relativo
entre 15°S y 30°S en verano, mientras que maximiza en latitudes medias
en invierno y primavera. Estas características están presentes también
en las otras ondas estacionarias consideradas.

En el caso de la QS2, se evidencia que a pesar de tener máximos de
\(r^2\) en la estratósfera al norte de 45°S, alcanza su máxima amplitud
al sur de esa latitud y en 200hPa en verano y en 30hPa en invierno. Su
actividad en la costa antártica se extiende en toda la tropósfera en
invierno (a pesar de que en \(r^2\) pierde importancia por encima de los
200hPa).

La región de amplitud máxima de la QS3, coincide aproximadamente con la
región de máximo \(r^2\) entre y otoño y primavera, aunque con menor
magnitud en superficie y menor extensión en toda la columna. En verano,
en cambio, exhibe un máximo de amplitud importante en latitudes medias
que no se asocia con valores máximos de \(r^2\).

La QS4 presenta un máximo de amplitud bien definido en la troposfera
media-alta sólo en verano, entre 15°S y 30°S.

\begin{landscape}\begin{figure}

{\centering \includegraphics[width=247mm,]{fig/tesis/ampl-ncep-1} 

}

\caption{Amplitud de Fourier (mgp) para números de onda 1 a 4 (NCEP).}\label{fig:ampl-ncep}
\end{figure}
\end{landscape}

\chapter{Onda 3}\label{onda-3}

En este capítulo se analiza la climatología observada de la QS3 a partir
de su reconstrucción mediante descomposición de Fourier. Se estudia su
amplitud y fase y se propone una división estacional alternativa a
partir del análisis de componentes principales. Además, se muestran la
estructura típica de la altura geopotencial y la función corriente
asociada a la QS3 a partir de regresiones con respecto a la amplitud de
la QS3.

\section{Características típicas}\label{caracteristicas-tipicas}

\begin{figure*}
\includegraphics[width=155mm, ]{fig/tesis/qs3-ncep-1} \caption{Z* media (mgp) reconstruida a partir de la QS3 en 300hPa (NCEP).}\label{fig:qs3-ncep}
\end{figure*}

En la \autoref{fig:qs3-ncep} se muestra el campo de Z* reconstruido sólo
a partir de la QS3 en 300hPa, que ilustra lo que sucede en todos los
niveles dado que su estructura es barotrópica equivalente (como se verá
en la \autoref{fig:qs3-ncep-corte}). De acuerdo con lo encontrado en la
amplitud de Fourier para la QS3 (\autoref{fig:ampl-ncep}), la amplitud
es máxima entre 60°S y 45°S con menor intensidad en primavera. En el
verano la fase es tal que una de las tres anomalías anticiclónicas se
ubica hacia el este del sur de Sudamérica. Se observa que existe un
corrimiento de la fase entre verano e invierno de poco más de 15° (algo
ya observado por \citet{Loon1972} y \citet{Mo1985}) anticipando que el
efecto de la QS3 sobre cada lugar pueda tener una componente estacional.

Además, las anomalías presentan una inclinación meridional que, como la
teoría de ondas de Rossby indica, está asociada a transportes
perturbados de cantidad de movimiento hacia el polo. La inclinación es
más importante en las estaciones de transición, pero menor en verano y
no detectable en el invierno.

\begin{figure*}
\includegraphics[width=155mm, ]{fig/tesis/qs3-ncep-corte-1} \caption{Corte zonal en 60°S de Z* (mgp) reconstruida a partir de la QS3 (NCEP).}\label{fig:qs3-ncep-corte}
\end{figure*}

La estructura vertical de las perturbaciones zonales de geopotencial
reconstruidas a partir de la QS3 se presenta en la
\autoref{fig:qs3-ncep-corte}. En todas las estaciones la mayor amplitud
se da en 300hPa, con el máximo desarrollo vertical en invierno, seguido
por el otoño y verano. Se destaca además la variación estacional de la
fase descripta previamente. Asimismo, se observa que existe una
variación estacional en la inclinación vertical de las perturbaciones.
En invierno y en primavera, los extremos de la onda presentan una ligera
inclinación hacia el oeste por debajo de los 300hPa, que no se detecta
en las otras estaciones.

\begin{figure*}
\includegraphics[width=155mm, ]{fig/tesis/qs3sd-ncep-1} \caption{Desvío estándar temporal de Z* (mgp) reconstruida a partir de la QS3. Se incluyen en negro los contornos de $\pm$ 20mgp de la amplitud de Fourier para indicar la posición de los centros de las perturbaciones (NCEP).}\label{fig:qs3sd-ncep}
\end{figure*}

En la \autoref{fig:qs3sd-ncep} se muestra el desvío estándar de Z*
reconstruidas a partir de la QS3. La variabilidad máxima se da entre los
centros de anomalías ciclónicas y anticiclónicas (marcados en la
\autoref{fig:qs3sd-ncep} con contornos negros), indicando que la
variabilidad del geopotencial asociada a la QS3 se asociaría
principalmente con defasajes que ocurren dentro de cada estación.

\section*{Wavelets}

\begin{figure*}
\includegraphics[width=155mm, ]{fig/tesis/wavelet-fourier-ncep-1} \caption{Amplitud (mgp) de la QS3 según wavelets (sombreados) y Fourier (contornos) (NCEP).}\label{fig:wavelet-fourier-ncep}
\end{figure*}

Como se describió en la \autoref{metodologia}, se utilizó la metodología
de wavelets como alternativa a la de Fourier para aislar la QS3. La
amplitud media obtenida mediante wavelets es prácticamente idéntica a la
amplitud obtenida con Fourier (\autoref{fig:wavelet-fourier-ncep}).

La metodología de wavelets también permite obtener información sobre la
variación meridional de la amplitud de la QS3, obteniendo así campos
horizontales de la amplitud de la QS3. Al igual que con Fourier, esto
puede hacerse promediando estacionalmente la amplitud de los campos
mensuales o calculando la amplitud de los campos estacionales. Los
valores de amplitud media zonal obtenidos por waveletes utilizando las
dos metodologías AM y MA son similares (no se muestra) pero sí hay
diferencias en las anomalías con respecto a la media zonal.

\begin{figure}

{\centering \includegraphics[width=155mm, ]{fig/tesis/waveletz-ncep-1} 

}

\caption{Anomalía zonal de la amplitud de la QS3 según wavelets (mpg) en 300hPa para el método AM y MA (\autoref{metodologia}) (NCEP).}\label{fig:waveletz-ncep}
\end{figure}

Las anomalías zonales de la amplitud de la QS3 en 300hPa, computadas con
wavelets se muestran en la \autoref{fig:waveletz-ncep} para los dos
métodos utilizados. Al norte de los 30°S, ambas metodologías muestran
resultados similares: en verano, otoño y primavera hay anomalías zonales
positivas sobre el Pacífico y negativas en el Índico, indicando que en
estas latitudes la QS3 tiene mayor amplitud en el hemisferio occidental
que en el oriental; en invierno, no hay anomalías zonales importantes en
esta región.

Al sur de 30°S las dos metodologías, en cambio, dan resultados
diferentes. Ambas presentan en todas las estaciones, una región elongada
de máxima amplitud en el sur del Índico centrado en 45°S que se extiende
hacia el este en latitudes más altas hasta alcanzar los 60°S
aproximadamente. La principal diferencia entre las metodologías,
entonces, es que la amplitud de la media estacional tiene los máximos
desvíos en 45°S, mientras que la media de la amplitud instantánea tiene
su máximo al sur de los 60°S.

\begin{figure*}
\includegraphics[width=155mm, ]{fig/tesis/estacionariedad-1} \caption{Estacionariedad de la QS3 (NCEP).}\label{fig:estacionariedad}
\end{figure*}

Estas diferencias en la \autoref{fig:waveletz-ncep} están relacionadas
con el grado de estacionariedad de las ondas. Se puede utilizar el
cociente entre la amplitud que resulta del método AM y la media de la
amplitud instantánea que resulta del método MA para definir un índice de
estacionariedad. Valores cercanos a 1, indican que el carácter
estacionario de las ondas sería el dominante, mientras que valores
cercanos a 0 implican que las ondas no tienen una localización
preferencial. La distribución horizontal de este índice de
estacionariedad se muestra en la \autoref{fig:estacionariedad}. La
máxima estacionariedad se observa en el verano tropical, pero esto se
debe a que la actividad de ondas estacionarias es casi nula.

Aproximadamente entre 30° y 60° se observa una banda de alta
estacionariedad, consistente con la región de mayor amplitud de QS3
(\autoref{fig:qs3-ncep}), que presenta un máximo al sur del Índico. Esto
sugiere que en esta región, la actividad de onda es particularmente
estacionaria, en contra de la sugerencia de \citet{Hobbs2010} de que el
principal patrón estacionario se da en el Pacífico Sur.

Estas observaciones destacan la utilidad de wavelets en el análisis de
ondas cuasiestacionarias. Mientras que el tratamiento con Fourier asume
a las ondas como una propiedad media de cada círculo de latitud,
wavelets permite reconocer su heterogeneidad meridional. Evaluando esta
heterogeneidad, sería posible distinguir entre campos donde una onda con
un determinado número de onda está presente en todo un círculo de
latitud de campos donde ésta está localizada en una región acotada.

Por otro lado, la no ortogonalidad de los wavelets complejiza la
interpretación de los resultados ya que no es posible la separación
estricta de un campo en modos oscilatorios con distinto número de onda.
El análisis de una QS específica, por lo tanto, está contaminado por la
actividad de otras QS con longitud de onda cercana.

El análisis de wavelets, en resumen, puede considerarse como
complementario al de Fourier agregando información de las asimetrías
zonales. Si embargo, dado que la variabilidad zonal es del orden de un
10\% de la amplitud media, en lo que sigue de la tesis se utilizará sólo
Fourier, dejando el análisis e interpretación de estudios de la QS3
utilizando wavelets para futuros trabajos.

\section{Amplitud}\label{amplitud}

Existen varios estadísticos que podrían utilizarse para representar la
amplitud de la QS3 en una región extendida, como la media, la máxima, la
moda o la mediana. En este caso, se estudió la posibilidad de
representarla con su media o su máxima, estimadas en la región
comprendida entre los 65°S y 40°S y entre 700hPa y 100hPa.

\begin{figure*}
\includegraphics[width=155mm, ]{fig/tesis/ampl-max-mean-1} \caption{Amplitud de la QS3 máxima y media (mgp) para 9 casos seleccionados.}\label{fig:ampl-max-mean}
\end{figure*}

\begin{figure*}
\includegraphics[width=155mm, ]{fig/tesis/ampl-max-mean-corte-1} \caption{Corte vertical de amplitud de la QS3 (mgp) para 9 casos seleccionados.}\label{fig:ampl-max-mean-corte}
\end{figure*}

\begin{figure*}
\includegraphics[width=155mm, ]{fig/tesis/ghz-ncep-select-1} \caption{Z* (mgp) en 300hPa con QS1 y QS2 eliminadas para 9 casos seleccionados.}\label{fig:ghz-ncep-select}
\end{figure*}

A modo exploratorio se seleccionaron 9 casos que representan distintas
características de la amplitud media y la máxima. Sus magnitudes
correspondientes se muestran en la \autoref{fig:ampl-max-mean}, el corte
vertical de la amplitud, en la \autoref{fig:ampl-max-mean-corte}, y los
campos de Z* (con las contribuciones de las QS1 y QS2 removidas como se
explicó en la \autoref{metodologia}) en la
\autoref{fig:ghz-ncep-select}.

Comparando el caso de mayo de 1997 con el de abril de 2012, ambos tienen
una amplitud media muy similar (\autoref{fig:ampl-max-mean}), pero la
amplitud máxima del primero es menor que la del segundo. El corte
vertical de la amplitud (\autoref{fig:ampl-max-mean-corte}) muestra que
la amplitud de la QS3 en abril de 2012 es grande y más localizada en la
troposfera alta que en el caso de mayo de 1997, para el cual la amplitud
es grande en la estratosfera. Esto se refleja en la comparación de las
perturbaciones de geopotencial (\autoref{fig:ghz-ncep-select}), donde la
QS3 en 300hPa se aprecia mucho más claramente en 2012 que en 1997.

Los casos de enero de 1985 y julio de 1988 también son similares en
amplitud media pero enero de 1985 tiene un valor mayor de amplitud
máxima. La estructura vertical de la amplitud muestra diferencias en la
extensión similares al par anterior de casos, con la amplitud en enero
de 1985 presentando un máximo bien definido en la tropósfera mientras
que en julio de 1988 alcanza valores importantes en niveles
estratosféricos. En este par de casos, los campos de Z* son muy
similares en cuanto a intensidad y claridad de la QS3. Ambos presentan
un tren de ondas que ocupan aproximadamente medio de círculo de latitud.
El tren de ondas de 1988 se ve algo más claro que el de 1985, a pesar de
que la amplitud máxima de 1985 es menor que la de 1988.

Los casos de noviembre de 1987 y enero de 2008 presentan características
diferentes a los anteriores. La amplitud máxima en ambos meses es
simliar, sin embargo en noviembre de 1987 la amplitud media es mayor que
en enero de 2008. En noviembre de 1987 la amplitud de la QS3 está más
extendida vertical y meridionalmente que en enero de 2008. Pero si se
observa el campo de Z* en 300hPa, si bien ambos casos tienen una QS3
clara, las anomalías en enero de 2008 son más intensas, especialmente en
la región del Índico.

Comparando los casos de diciembre de 1990 y septiembre de 2000, puede
verse que ambas medidas de amplitud son mayores en diciembre de 1990 que
en septiembre de 2000 y que ambos meses tienen un desarrollo vertical
similar de la amplitud de la QS3. El campo de Z* en 300hPa muestra en
diciembre de 1990 una estructura de la QS3 más zonal y coherente a lo
largo de todas las longitudes. Sin embargo, en septiembre de 2000, se
destaca principalmente un tren de ondas, asociado con anomalías más
intensas, por lo que su efecto en la circulación local puede ser mayor
que las de diciembre de 1990.

Más extrema aún es la diferencia entre los casos de septiembre de 2000 y
octubre de 2003. Ambos tienen métricas de amplitud similares tanto en
media como en máxima y una extensión vertical de la QS3 similar. No
obstante, la QS3 en 300hPa de octubre de 2003 es apenas distinguible en
comparación con septiembre de 2000.

Estos casos ilustran, entonces, la diversidad de situaciones en que
puede desarrollar la QS3 y las limitaciones que un estudio
dinámico-estadístico como el que sigue, puede tener. Algunos problemas
son inherentes a cualquier intento de representar una estructura con
variabilidad espacial a partir de un sólo número o indicador y otros
están relacionados con la limitación de la descomposición de Fourier que
trata toda onda como una onda planetaria con igual estructura a lo largo
de todo un círculo de latitud.

\begin{figure*}
\includegraphics[width=155mm, ]{fig/tesis/cor-mean-max-1} \caption{Amplitud máxima y media de la QS3 para cada mes del período 1985-2015 (mgp). La línea sólida representa la regresión lineal entre ambas variables ($r^2 = 0.94 $) y la línea punteada es la línea amplitud máxima = amplitud media.}\label{fig:cor-mean-max}
\end{figure*}

Cabe aclarar además que estos 9 casos fueron seleccionados
específicamente para ilustrar estas limitaciones y no son necesariamente
representativos de la totalidad de casos posibles. Como se muestra en la
\autoref{fig:cor-mean-max}, la amplitud media y la máxima tienen una
relación lineal con un coeficiente alto de correlación asociado
(\(r^2>0,9\)). Debido a esto, a fines estadísticos la elección de una u
otra métrica no tendría, en promedio, una influencia importante (aunque
podría tenerla en casos individuales, como lo ilustraron los casos
previos).

A partir de esta discusión, se definió un índice de la QS3 (\(A_3\))
como el promedio espacial de la amplitud (computada a partir de cada
media mensual) entre 65°S y 40°S y entre 700hPa y 100hPa.

El ciclo anual de \(A_3\) se muestra en la Figura~\ref{fig:ampl-ts1} a
través de boxplots de los correspondientes valores mensuales
individuales. Los valores mayores de amplitud se dan en agosto y los
mínimos en diciembre, y la variabilidad sigue un ciclo similar. Este
ciclo anual es consistente con trabajos previos
\citep{Loon1972, Karoly1985, Raphael2004}. Se nota que \(A_3\) no es
mínimo en los meses de primavera, en contraste con lo observado en el
análisis climatológico (\autoref{fig:ampl-ncep}) y los campos
reconstruidos (\autoref{fig:qs3-ncep}). Esta aparente contradicción
surge, como se verá más adelante, en la falta de considerar también la
fase de la QS3 (\autoref{fase}).

\begin{figure*}
\subfloat[Boxplot para cada mes. La línea central representa la mediana, la caja se extiende desde el percentil 25\% hasta el percentil 75\% y las líneas se extienden desde los extremos de la caja hasta 1,5 veces el intervalo intercuartil. Los puntos son los valores mensuales individuales.\label{fig:ampl-ts1}]{\includegraphics[width=155mm, ]{fig/tesis/ampl-ts-1} }\newline\subfloat[Serie temporal. Las líneas horizontales representan la media anual, en rojo (azul) cuando ésta es mayor (menor) que la media de todo el período marcada con línea punteada.\label{fig:ampl-ts2}]{\includegraphics[width=155mm, ]{fig/tesis/ampl-ts-2} }\caption{Índice $A_3$ (mgp).}\label{fig:ampl-ts}
\end{figure*}

La serie temporal de \(A_3\) se muestra en la Figura~\ref{fig:ampl-ts2}.
La serie exhibe tanto variaciones interanuales como subanuales de la
amplitud. Además, se observan series de años con anomalías positivas
consistentemente seguidos por anomalías negativas (1985-1990, 1992-1996
y 1999-2005) y otras con persistencia de anomalías positivas o negativas
(2005-2009 y 2012-2015). No hay evidencia visual de periodicidades ni de
una tendencia secular.

Un análisis de wavelets de la serie de anomalías temporales de \(A_3\)
(definidas a partir de la remoción del ciclo anual) presenta picos
alrededor de 3, 12 y 24 meses, aunque ninguno es estadísticamente
significativo a 95\% de confianza (no se muestra). \citet{Raphael2004},
también utilizando wavelets, pero en el período 1958-2001, encontró que
su índice de actividad de la QS3 (discutido en la \autoref{metodologia})
muestra variabilidad importante entre 4 y 12 años, algo que no se
encontró en este estudio, quizás porque cada índice captura
características diferentes o porque los períodos de los datos analizados
son distintos.

\section{Fase}\label{fase}

Resulta evidente que además de la amplitud, las ondas planetarias se
caracterizan por su fase. Los resultados de trabajos previos incluidos
en la Introducción (como por ejemplo los de \citet{Trenberth1985}) y así
como los de las secciones anteriores muestran que la fase de la QS3
puede presentar cambios importantes. En este trabajo se definió el
índice \(F_3\) como el promedio espacial de la fase (computada a partir
de cada media mensual) entre 65°S y 40°S y entre 700hPa y 100hPa. La
fase se computó de manera que represente la longitud del máximo de
perturbación de Z* explicada por la QS3 entre 120°O y 0° de longitud.

La \autoref{fig:fase-boxplot} muestra, para cada mes, en rojo los
valores de \(F_3\) y en azul \(F_3\pm 60^{\circ}\) (representando la
longitud del \emph{mínimo} de Z* explicado por la QS3) para los 20 años
con un valor de \(A_3\) más extremo. Se incluye el promedio y el rango
delimitado por \(\pm\) 1 desvío estándar. El mapa se muestra para
referencia. Si bien la posición de los puntos en la dirección horizontal
tiene relación con la longitud, en el eje vertical no tiene relación con
la latitud y sí en cambio con los meses del año.

\begin{landscape}\begin{figure}

{\centering \includegraphics[width=247mm,]{fig/tesis/fase-boxplot-1} 

}

\caption{$F_3$ (grados) para cada mes del año a partir de los 20 años con un valor de $A_3$ más extremo y el rango definido por $\pm$ 1 desvío estándar  (puntos negros y barras negras). En rojo y azul se identifica respectivamente la localización del máximo y el mínimo de perturbación de Z para cada año individual}\label{fig:fase-boxplot}
\end{figure}
\end{landscape}

El ciclo anual de \(F_3\) que se obtiene de esta manera es coherente con
el descripto en la \autoref{fig:qs3-ncep}. La fase se centra en promedio
en 55°O en enero y en 90°O en junio con ubicaciones intermedias en los
meses de transición. En particular, el continente queda al este del
máximo de anomalía de geopotencial en invierno lo que, por balance
geostrófico, implicaría anomalías positivas del viento meridional, que
favorecerían los vientos del sur. En verano, en cambio, se da la
situación contraria. Esto confirma que el efecto de la QS3 sobre
Sudamérica depende crucialmente de su fase.

La variabilidad interanual de la fase para cada mes es considerable y de
una magnitud comparable a la amplitud del ciclo anual. En particular, es
notorio el aumento en la variabilidad de la fase en los meses de
primavera, al punto de que en noviembre la fase prácticamente no tiene
una posición predilecta.

La gran variabilidad presente durante los meses de primavera, en
comparación con el resto del año, explica por qué en los campos medios
la QS3 aparece débil en esa estación, a pesar de que su amplitud mensual
no es pequeña. Como se mencionó en la \autoref{metodologia}, promediar
campos que están defasados en entre 1/3 y 2/3 de longitud de onda (entre
40° y 80° en el caso de la QS3) interfiere destructivamente entre ellos,
eliminando la señal en los campos medios. En primavera, más del 18\% de
los meses tienen algún grado de interferencia destructiva (estimada como
se describió en la \autoref{metodologia}), porcentaje mayor a, por
ejemplo, el encontrado del 8\% en verano.

Observando la distribución de los centros de anomalías ciclónicas y
anticiclónicas (puntos azules y rojos, respectivamente) se nota que, a
pesar de que Sudamérica está principalmente afectada por centros
anticiclónicos asociados a la QS3, existe un número no despreciable de
años en los que se observa un centro ciclónico sobre el continente. En
particular, sumando noviembre y diciembre se pueden observar 9 años
donde esto sucedió.

\section{Estaciones}\label{estaciones}

En las secciones anteriores se mostraron campos medios estacionales
utilizando la definición tradicional de las estaciones climatológicas
(verano = DEF, otoño = MAM, invierno = JJA, primavera = SON). Sin
embargo, como éstas son definidas a partir del ciclo anual de
temperatura en latitudes medias, no constituyen necesariamente el mejor
agrupamiento de los datos para otras variables u otras latitudes (por
ejemplo, la Antártida \citep{King1997}).

Una metodología muy utilizada para la clasificación de campos es el
análisis de componentes principales (PCA). Como se explicó en la
\autoref{metodologia}, se aplicó esta metodología a los campos mensuales
de Z* reconstruidos a partir de la QS3 pesados por la raíz cuadrada del
coseno de la latitud para caracterizar más objetivamente su
variabilidad.

La \autoref{eoftabla} muestra la varianza explicada de cada componente
principal obtenida a partir de los campos reconstruidos de QS3. Las
primeras dos componentes explican en conjunto más del 80\% de la
varianza, explicando cada una parte similar de la varianza (lo que
indica que se trata de autovalores degenerados). Sabiendo, además, que
los campos de QS3 prácticamente sólo tienen dos grados de libertad
(amplitud y fase), la elección de seleccionar las dos primera
componentes resulta natural además de justificada por la
\autoref{eoftabla}

\begin{longtable}[]{@{}lrrrrr@{}}
\caption{Varianza explicada por las 5 primeras componentes principales
de los campos de QS3 reconstruidos.\label{eoftabla}}\tabularnewline
\toprule
& PC 1 & PC 2 & PC 3 & PC 4 & PC 5\tabularnewline
\midrule
\endfirsthead
\toprule
& PC 1 & PC 2 & PC 3 & PC 4 & PC 5\tabularnewline
\midrule
\endhead
\(r^2\) & 0,436 & 0,397 & 0,054 & 0,038 & 0,026\tabularnewline
\bottomrule
\end{longtable}

\begin{figure*}
\includegraphics[width=155mm, ]{fig/tesis/eof-field-1} \caption{Primeras dos componentes principales del campo de Z* reconstruido a partir de la QS3.}\label{fig:eof-field}
\end{figure*}

Las dos primeras componentes principales del campo de QS3
(\autoref{fig:eof-field}) se encuentran en cuadratura y su combinación
lineal resulta en otra QS3 cuya fase depende de la amplitud relativa de
cada componente. En verano predomina la PC1, mientras que en invierno
predomina la PC2 como se muestra en la \autoref{fig:pc1-pc2}.

Se observa que enero, febrero y marzo tienen preponderancia de la PC1 y
casi nulo de la PC2 por lo que podrían agruparse. Abril, mayo, agosto,
septiembre y octubre tienen una mezcla similar de componentes, pero es
conveniente separar los dos primeros para respetar la progresión
temporal. Junio y julio tienen un comportamiento diferente al de los
demás meses con gran magnitud de PC2. Noviembre y diciembre aparecen
como \emph{outliers} en este diagrama debido a que su mayor variabilidad
(como se notó en la \autoref{fig:fase-boxplot}) hace que no predomine
ninguna componente principal. Podrían ser clasificarlos juntos como
meses de ``no estacionariedad'' indicando que se trata de una época del
año donde la QS3 no está presente.

\begin{figure*}
\includegraphics[width=155mm, ]{fig/tesis/pc1-pc2-1} \caption{Valor medio de las dos primeras componentes principales del campo de Z* reconstruido a partir de la QS3 para cada mes. Las líneas unen cada mes siguiendo el orden anual y los colores separan a las 5 "estaciones" definidas en el texto.}\label{fig:pc1-pc2}
\end{figure*}

La descripción de la QS3 basada en esta nueva descomposición estacional
se presenta en la \autoref{fig:qs3-qsseason-ncep}, donde se muestra el
campo de Z* reconstruido sólo a partir de la QS3 en 300hPa. En la
\autoref{fig:qs3-qsseason-ncep-corte} que muestra el corte
vertical-zonal correspondiente a 52,5°S. Comparando con las Figuras
\ref{fig:qs3-ncep} y \ref{fig:qs3-ncep-corte} de la
\autoref{caracteristicas-tipicas} se ve que en las estaciones
alternativas los campos de las estaciones de transición (AM y ASO) son
más similares entre, sí tanto en el campo horizontal como en el corte
meridional. Como era de esperarse, el bimestre ND prácticamente carece
de valores medios significativos asociados con la QS3.

\begin{figure*}
\includegraphics[width=155mm, ]{fig/tesis/qs3-qsseason-ncep-1} \caption{Z* media reconstruida a partir de la QS3 en 300hPa (mgp) según las estaciones definidas en el texto.}\label{fig:qs3-qsseason-ncep}
\end{figure*}

\begin{figure*}
\includegraphics[width=155mm, ]{fig/tesis/qs3-qsseason-ncep-corte-1} \caption{Corte en 52,5°S de Z* media reconstruida a partir de la QS3 en 300hPa (mgp) según las estaciones definidas en el texto.}\label{fig:qs3-qsseason-ncep-corte}
\end{figure*}

Además, la inclinación vertical de las perturbaciones que en la
\autoref{fig:qs3-ncep-corte} aparece sólo levemente en invierno se
observa con más claridad con la nueva descripción estacional
(\autoref{fig:qs3-qsseason-ncep-corte}). En el bimestre JJ la
inclinación hacia el oeste es importante en la troposfera, estando
presente también en AM y ASO en menor medida.

La teoría de ondas de Rossby indica que una inclinación hacia el oeste
está asociada con transporte de calor hacia el polo por las
perturbaciones y propagación vertical de las mismas \citep{James}. Esto
podría explicar la variación estacional de la extensión vertical de la
QS3, la cual tendría condiciones favorables para desarrollarse en la
estratósfera alta en JJ, cuando la inclinación es máxima y quedaría
atrapada en niveles más bajos en verano cuando la inclinación es mínima.

\begin{figure*}
\includegraphics[width=155mm, ]{fig/tesis/lag-cor-1} \caption{Correlación lageada para cada mes con los 12 siguientes. Valores a la derecha de la línea indican correlación entre meses del mismo año y a la izquierda, entre un mes con el del año siguiente. Los colores indican el valor de la correlación. }\label{fig:lag-cor}
\end{figure*}

De manera de explorar la persistencia de la actividad de la QS3, se
presenta la \autoref{fig:lag-cor}, que muestra la correlación desfasada
del campo de QS3 correspondiente a cada mes con todos los demás. Es
decir, el valor de enero con el de febrero siguiente representa la
correlación entre los campos de QS3 de todos los eneros con los campos
de todos los febreros que les siguen. La línea escalonada marca la
separación del año de manera que un número a la izquierda de ésta
implica correlación de ese mes con el mismo mes del año siguiente. Las
correlaciones justo a la izquierda de la línea escalonada son positivas
y relativamente altas para casi todos los meses salvo noviembre,
diciembre y agosto. Esto implica que el comportamiento de la QS3 en
estos meses son similares de un año a otro. Noviembre y diciembre
también presentan bajas correlaciones en general con los demás meses, lo
que es coherente con los resultados de las secciones anteriores, es
decir, al ser meses con actividad de la onda 3 poco estacionaria, sus
campos de QS3 no son consistentemente similares con ningún otro mes.
Esta interpretación no parece posible para agosto, ya que su
variabilidad no es particularmente alta.

Los valores de un mes a la derecha de la línea escalonada son en
generalmente altos indicando buena concordancia entre los campos de un
mes y el siguiente. Para esto nuevamente las excepciones son noviembre,
diciembre y julio. Las correlaciones entre meses corridos 6 meses son
bajas para los meses de verano e invierno y medias para los meses de
transición. Es decir, los meses de verano son muy distintos de los de
invierno, mientras que los de transición son medianamente parecidos a
todos. Esta es una consecuencia del ciclo anual de la fase
(\autoref{fig:fase-boxplot}).

Se observa que los máximos valores de correlación se dan entre enero,
febrero y marzo, indicando un nivel de persistencia relativamente alto.
Esto puede deberse a la baja variabilidad de la fase presente en estos
meses o también a la influencia de los eventos ENSO, cuyos precursores
podrían estar relacionados con una onda 3 durante el verano
\citep{Qin2017}.

El uso de componentes principales para el análisis de una onda que
cambia de fase es similar a la metodología utilizada para el monitoreo
de la MJO \citep{Wheeler2004}, por lo que sería posible su utilización
como indicador de la actividad de la QS3 distinto a \(A_3\). La
exploración de dicho indicador está por fuera del objetivo de este
trabajo.

\section{\texorpdfstring{\(r^2\)}{r\^{}2}}\label{r2}

En el \autoref{ondas-quasiestacionarias} se discutió la relevancia del
parámetro \(r^2\) como indicador de las características de la QS3. En
particular, se mostró que la estructura vertical de la varianza
explicada por la QS3 para cada estación (\autoref{fig:r2-ncep}) se
caracteriza por máximos valores cerca de 45°S entre la superficie y
200hPa y un ciclo anual con máximo en invierno y mínimo en verano. En
esta sección se explora su distribución horizontal para la QS3. Para
esto, se definió el \(r^2\) como la correlación cuadrada entre el valor
de Z* y el de QS3 correspondiente para cada punto de grilla y cada mes.

\begin{figure*}
\includegraphics[width=155mm, ]{fig/tesis/cor-campo-1} \caption{$r^2$ entre Z* en 300hPa y QS3 (sombrados). Se incluyen en  los contornos de $\pm$ 20mgp de la amplitud de Fourier para indicar la posición de los centros de las perturbaciones, en rojo los centros anticiclónicos y en azul, los ciclónicos (NCEP).}\label{fig:cor-campo}
\end{figure*}

Los campos horizontales de \(r^2\) en 300hPa se muestran en la
\autoref{fig:cor-campo}. En las cuatro estaciones la QS3 explica la
mayor parte de la varianza en el hemisferio occidental entre 60°S y
45°S. Lejos de ser homogéneos, los campos muestran tres máximos
localizados con cierta persistencia durante el año. El primero, al sur
del Índico, está presente en verano y otoño en 60°E que se encuentra más
hacia el este en invierno y primavera. Algo similar sucede con el
segundo máximo ubicado al sur del Pacífico, que se encuentra en 180° en
otoño y en 120°O en primavera. Se observa un tercer máximo en el
Atlántico sur cuya ubicación varía longitudinalmente poco. En verano y
otoño se distingue un máximo en latitudes bajas en el pacífico central.

Si se compara la posición de los máximos de \(r^2\) con los centros de
QS3, en verano, el máximo del Índico, por ejemplo, coincide con un
centro anómalamente anticiclónico en verano. No obstante, el máximo de
\(r^2\) de otoño sobre esa cuenca se encuentra entre dos centros de la
onda y lo mismo pasa con el máximo del Atlántico. Éste último coincide
con un centro de anomalía anticiclónica en verano, aunque en invierno
está más cerca de uno ciclónico. En suma, no parece haber una asociación
entre los máximos de \(r^2\) y la ubicación de los centros de la QS3.

Se realizó el mismo análisis en base a las estaciones definidas según
las características de la QS3 (\autoref{estaciones}) y los resultados
generales no se modifican. Las principales diferencias son un
debilitamiento de los máximos de \(r^2\) durante EFM y una fuerte
intensificación del máximo del Atlántico durante ND (no se muestra).

\section{Regresiones}\label{regresiones}

Los campos de Z* explicados por la QS3 descripta por el método de
Fourier proporcionan una descripción inicial pero limitada de la
influencia de esta onda sobre las anomalías de circulación ya que no
permiten describir completamente el estado típico de la atmósfera cuando
la QS3 está activa. En esta sección se exploran brevemente las anomalías
típicas de esas variables explicadas por la actividad de la QS3, a
través del cómputo de mapas de regresión lineal entre \(A_3\)
(\autoref{amplitud}) y diferentes variables. En esta metodología, la
serie de \(A_3\) está estandarizada de manera que los valores de
regresión representan la porción de anomalía de la variable en cuestión
asociado con un cambio de la amplitud de la QS3 de 1 desvío estándar de
magnitud. En ese sentido los valores de regresión tienen las unidades de
la variable regresionada. Para facilitar la interpretación de los
resultados, se los describe asumiendo que las regresiones corresponden a
un cambio positivo en la amplitud de la QS3 de 1 desvío estándar.

\section*{Geopotencial}

\begin{landscape}\begin{figure}

{\centering \includegraphics[width=247mm,]{fig/tesis/regr-gh-ncep-1} 

}

\caption{Regresión de Z en 300hPa con $A_3$ estandarizado (mgp).}\label{fig:regr-gh-ncep}
\end{figure}
\end{landscape}

En la \autoref{fig:regr-gh-ncep} se muestra la regresión del campo de Z
en 300hPa con \(A_3\) para los doce meses del año. Si bien todos
presentan evidencias visibles de la actividad de la QS3, presentan
también mucha heterogeneidad en las características observadas en los
distintos meses.

En enero, se observa un patrón hemisférico de onda 3, con un centro
anómalamente negativo importante en el Pacífico sudeste, embebido además
en un tren de ondas extendido entre el este de Nueva Zelanda, y
Sudamérica. En febrero, la estructura se mantiene similar, pero con
menor evidencia de propagación meridional, y el centro anticiclónico en
Sudamérica intensificado y un debilitamiento de aquel ubicado al sur del
Índico. En marzo, la estructura de onda 3 exhibe mínima variación
meridional y ocupa todo el círculo de latitud entre 60°S y 45°S. Abril
es similar a marzo, pero sin centros ciclónicos significativos. En mayo,
se distingue nuevamente un tren de ondas con propagación meridional,
emanando desde Australia, pero con evidencias de reflexión en el Mar de
Weddell. Junio, al igual que marzo, presenta una estructura de QS3
zonal, pero menos definida y, consistente con el corrimiento de la fase
observado (\autoref{fase}), con los centros ciclónicos y anticiclónicos
desplazados hacia el oeste.

En julio, el patrón de QS3 se encuentra superpuesto con un fuerte patrón
anular, con anomalías positiva de geopotencial en latitudes polares y
negativas en latitudes medias. En la \autoref{fig:regr-gh-polar} se
muestran los mismos campos para julio y diciembre que en la
\autoref{fig:regr-gh-ncep} pero en proyección estereográfica polar. En
esta proyección es fácil identificar un patrón anular similar en julio a
una fase negativa del SAM y a una fase positiva del SAM en diciembre. En
efecto, la correlación entre el índice SAM y \(A_3\) en julio es -0,57
(p-valor \(\sim 9\times 10^{-4}\)) y en diciembre es 0,41 (p-valor
\(\sim 0,0225\))\footnote{Correlación de Pearson, tests a dos colas.}.
Estos dos meses son los únicos con una relación significativa con el SAM
al nivel del 95\% de confianza.

\begin{figure*}
\includegraphics[width=155mm, ]{fig/tesis/regr-gh-polar-1} \caption{Idem \autoref{fig:regr-gh-ncep}, pero en coordenadas estereográficas polares para julio y diciembre.}\label{fig:regr-gh-polar}
\end{figure*}

La \autoref{fig:regr-gh-ncep} muestra que en agosto el campo de
regresión tiene una estructura de QS3 zonal salvo en la región de
Sudamérica, donde el centro anticiclónico se encuentra a mayor latitud,
sobre la Península Antártica. En septiembre se observa, embebido sobre
el patrón de onda 3, un tren de ondas sobre el Pacífico con importante
propagación meridional, similar al de otros meses (como enero, por
ejemplo). En octubre, el patrón de QS3 se ve claramente en los centros
de anomalías positivas, pero no tanto en los negativos. En este mes, hay
evidencias de propagación meridional de trenes de onda. En particular,
se distingue un tren de ondas extendido entre el Pacífico central y el
Atlántico, con fase tal que exhibe un centro negativo sobre Sudamérica.
El campo de regresión en noviembre es muy similar al de abril; con
evidencias de un tren de ondas extendido entre el sur de Australia y
Sudamérica. En diciembre el patrón es similar al de enero pero
modificado por el patrón anular similar a la fase positiva del SAM,
discutido previamente.

\section*{Función Corriente}

\begin{landscape}\begin{figure}

{\centering \includegraphics[width=247mm,]{fig/tesis/regr-psi-ncep-1} 

}

\caption{Regresión de la anomalía zonal de $\psi$ en $\sigma = 0,2101$ con $A_3$ estandarizado ($1\times10^{-6}m^2/s$) y flujos de actividad de onda calculados a partir de la misma.}\label{fig:regr-psi-ncep}
\end{figure}
\end{landscape}

Como se mostró previamente, la función corriente (\(\psi\)) resulta una
mejor variable que Z para describir anomalías de circulación tanto sobre
latitudes ecuatoriales como polares. Asimismo, los mapas de regresión de
Z explicadas por la QS3 muestran en varios meses del año evidencias de
trenes de onda extendidos entre zonas tropicales y extratropicales. En
este sentido, esta subsección describe las características de la
regresión entre las anomalías zonales de \(\psi\) con \(A_3\). Se
incluye además los flujos de actividad de onda calculados a partir de
tal regresión. El estudio de estos flujos permite confirmar (o no) que
una determinada alternancia de centros de diferente signo está
relacionada (o nó) con un tren de onda.

La \autoref{fig:regr-psi-ncep} muestra evidencias de propagación
meridional de ondas de Rossby desde la zona ecuatorial y atravesando el
Pacífico Sur hasta Sudamérica en varios meses del año aunque con
diferentes características.

En enero y febrero, trenes de onda parecen emanar desde la zona
ecuatorial ubicada al noreste de Nueva Zelanda, aunque con diferente
fase y reflejadas a diferentes longitudes. En marzo, los flujos
confirman la dominancia de la dispersión zonal de energía de ondas
Rossby entre el Indico y el Pacífico, mientras que, en abril, las
señales son muy débiles. En mayo, se observa un tren intenso desde el
norte de Australia hacia el sur de Sudamérica. En junio la actividad de
onda que emana desde Australia lo hace principalmente a lo largo de los
subtrópicos para luego, en el Pacífico sudeste, dispersarse hacia el sur
e influenciar el sur de Sudamérica. En agosto y septiembre, la actividad
de onda que influye el sur de Sudamérica emana principalmente de la
vecindad de Nueva Zelanda. En octubre, domina el flujo de actividad de
onda que emana desde el norte de Nueva Zelanda hacia el Pacífico
ecuatorial central para luego dispersar hacia el sudeste hasta
Sudamérica y el Atlántico Sur. En noviembre, se distingue una dispersión
de energía desde el sur de Australia, hacia el sur, que se observa más
fuerte y coherente en diciembre influenciando Sudamérica.

Cabe mencionar que se hicieron también análisis de regresión entre las
anomalías de la temperatura de la superficie del mar y \(A_3\) con el
fin de explorar en qué medida los trenes de onda descriptos previamente
están relacionados con variaciones en la temperatura de superficie del
océano, en especial de los océanos tropicales, pero no se encontraron
relaciones significativas. Las limitaciones que tiene la descripción de
las características de la QS3 solo a partir de la amplitud, y también la
imposibilidad de los análisis con observaciones de poder separar
limpiamente los factores asociados con todas las posibles causas,
motivaron la realización de los experimentos con el modelo SPEEDY que se
discuten en el capítulo siguiente.

\chapter{Simulaciones con el modelo
SPEEDY}\label{simulaciones-con-el-modelo-speedy}

El modelo SPEEDY fue utilizado para explorar las posibles causas que
expliquen el desarrollo de un patrón de onda 3 en la circulación del HS,
más allá de las causas relacionadas con las vacilaciones del flujo de
los oestes. En este capítulo se presentan y discuten los resultados de
la corrida Control, así como de los experimentos de sensibilidad
realizados.

\section{Validación}\label{validacion}

Se realizó primero una corrida Control utilizando el modelo SPEEDY. La
misma cubrió el período 1985-2014, con los modelos de suelo y hielo
activados y considerando como forzante las medias mensuales de SST
provenientes de la base de datos HadISST \citep{Rayner2003}. Se utilizó
una atmósfera en reposo como condición inicial observando que el tiempo
de \emph{spin-up} es menor que un mes, por lo que se utilizaron todos
los meses de simulación en el análisis.

Para la validación de la corrida Control se utilizaron los datos del
NCEP, descriptos en la \autoref{datos-y-modelo} y utilizados en los
Capítulos \ref{climatologia-observada} y \ref{onda-3}. La descripción se
concentra en mostrar los resultados para diferentes variables en los
niveles de 200hPa y 500hPa, como representativos de los niveles altos y
medios de la troposfera respectivamente. Las conclusiones no cambian
substancialmente en el resto de los niveles.

Como se mencionó previamente (\autoref{datos-y-modelo}), la
representación de la estratósfera en SPEEDY es muy limitada debido a que
el único nivel estratosférico disponible (30hPa) es la tapa del modelo,
funcionando como ``esponja'' que evita la propagación de ondas de
gravedad. Teniendo en cuenta, los resultados de trabajos previos y del
\autoref{climatologia-observada}, donde se presentan evidencias de
actividad de la QS3 en la estratosfera invernal, se asume que esto puede
limitar la calidad de las simulaciones durante esa época del año.

\subsection{Altura Geopotencial}\label{altura-geopotencial-1}

\begin{landscape}\begin{figure}

{\centering \includegraphics[width=247mm,]{fig/tesis/ghz-sp-nc-1} 

}

\caption{Z* (mgp) (SPEEDY sombreado, NCEP contornos).}\label{fig:ghz-sp-nc}
\end{figure}
\end{landscape}

\begin{figure*}
\includegraphics[width=155mm, ]{fig/tesis/ghz-sp-nc-corte60-1} \caption{Corte zonal de Z* (mgp) en 60°S (SPEEDY sombreado, NCEP contornos).}\label{fig:ghz-sp-nc-corte60}
\end{figure*}

La validación del campo medio de altura geopotencial (no se muestra)
confirma que está bien representado por el modelo. La comparación de las
anomalías zonales del campo medio representadas por SPEEDY se muestra en
la \autoref{fig:ghz-sp-nc} (campo horizontal) y en la
\autoref{fig:ghz-sp-nc-corte60} (corte zonal en 60°S), donde en
sombreado se muestra el campo de SPEEDY y en contornos el
correspondiente a NCEP (convención que se mantendrá en el resto de las
figuras de validación). SPEEDY representa correctamente la estructura
aproximadamente barotrópica equivalente de las anomalías zonales y la
predominancia de la señal asociada con la QS1. Sin embargo, las
simulaciones presentan errores en la ubicación, intensidad e inclinación
de las perturbaciones.

En invierno, las anomalías simuladas son de magnitud importante como las
observadas, aunque el modelo no representa la inclinación observada de
las mismas hacia el oeste con la altura
(\autoref{fig:ghz-sp-nc-corte60}). Este error es coherente con la falta
de representación de la estratosfera en el modelo y que los estudios
previos relacionan la inclinación hacia el oeste con propagación
vertical de ondas de Rossby desde la troposfera hacia la estratosfera.
Sin embargo el hecho que las anomalías en la troposfera estén
razonablemente simuladas (aunque algo más intensas que las observadas),
da evidencias la importancia de los procesos troposféricos en mantener
estas anomalías. Las anomalías simuladas en otoño y primavera presentan
características similares a las invernales. En cambio, en verano el
patrón es bastante más débil que el observado. Se destaca además que si
bien en líneas generales la ubicación de los máximos y mínimos de Z* en
SPEEDY es aproximadamente la correcta, no logra capturar parte de la
estructura fina. Las anomalías simuladas en 500hPa en otoño, invierno y
primavera, presentan un sólo máximo en 120°O a pesar de que NCEP muestra
dos máximos, uno que se mantiene en la línea de fecha durante todo el
año y otro entre 120°O y 60°O según la época del año.

\begin{landscape}\begin{figure}

{\centering \includegraphics[width=247mm,]{fig/tesis/ghz-dif-sp-nc-1} 

}

\caption{Diferencia de Z* (mgp) entre SPEEDY y NCEP.}\label{fig:ghz-dif-sp-nc}
\end{figure}
\end{landscape}

En la \autoref{fig:ghz-dif-sp-nc} se muestra la diferencia entre el
campo de Z* de NCEP y SPEEDY. En verano y otoño la principal diferencia
radica en que NCEP muestra una alta más intensa de la QS1 al norte de
60°S. En invierno y primavera se observa un tren de ondas con
propagación meridional que une el Índico con el Atlántico con número de
onda planetaria 3. Un tren de ondas similar fue identificado en las
observaciones en la \autoref{viento-meridional}. Su aparición al hacer
la resta NCEP - SPEEDY indica que el mismo no está presente en la
corrida Control.

\subsection{Temperatura}\label{temperatura-1}

\begin{landscape}\begin{figure}

{\centering \includegraphics[width=247mm,]{fig/tesis/t-nc-sp-1} 

}

\caption{Campo medio de temperatura (°C) (SPEEDY sombreado, NCEP contornos).}\label{fig:t-nc-sp}
\end{figure}
\end{landscape}

En la \autoref{fig:t-nc-sp} se muestra el campo medio de temperatura de
SPEEDY y NCEP. En 850hPa y 500hPa, ambos campos son muy similares. En
niveles altos (200hPa) la simulación control difiere considerablemente
de las observaciones. En verano y en otoño, SPEEDY muestra un gradiente
meridional mucho más importante que NCEP y en invierno y primavera el
gradiente máximo se da entre 30°S y 45°S para SPEEDY, y en 60°S en NCEP.

La \autoref{fig:tz-sp-nc} presenta las anomalías zonales de temperatura
simuladas y observadas. La estructura de ambas coincide en niveles
bajos, donde la influencia superficial es importante, pero son
diferentes en altura. En 500hPa, las anomalías Antárticas de SPEEDY
coinciden aproximadamente en ubicación con las de NCEP aunque son
ligeramente más débiles en invierno y primavera. En 60°S, en verano hay
buena coincidencia, pero entre otoño e invierno la QS1 observada en esas
latitudes virtualmente desaparece en SPEEDY mientras que en primavera
vuelve a crecer incluso con mayor intensidad que en NCEP. En 200hPa,
SPEEDY carece casi totalmente de anomalías zonales significativas al sur
de los 45°S durante todo el año a diferencia de NCEP, que muestra una
estructura de QS1 bien definida con máxima amplitud en primavera.
Nuevamente, esto puede deberse a la falta de representación de la
estratosfera en el modelo. Al norte de esa latitud las anomalías de
SPEEDY coinciden mejor con NCEP.

\begin{landscape}\begin{figure}

{\centering \includegraphics[width=247mm,]{fig/tesis/tz-sp-nc-1} 

}

\caption{Anomalía zonal de temperatura (°C) (SPEEDY sombreado, NCEP contornos)}\label{fig:tz-sp-nc}
\end{figure}
\end{landscape}

\subsection{Viento zonal}\label{viento-zonal-1}

\begin{figure*}
\includegraphics[width=155mm, ]{fig/tesis/u-sp-nc-corte-1} \caption{Media zonal del viento zonal (m/s) (SPEEDY sombreado, NCEP contornos).}\label{fig:u-sp-nc-corte}
\end{figure*}

En la \autoref{fig:u-sp-nc-corte} se muestra el viento zonal medio para
SPEEDY y NCEP. El modelo es capaz de representar varias de las
características observadas más importantes, como ser la existencia de
regiones de vientos del oeste intensos en latitudes subtropicales y
subpolares, así como sus variaciones estacionales. Sin embargo, el jet
simulado en verano se encuentra más al norte y ligeramente más elevado,
así como los jets subtropicales son simulados considerablemente más
intensos en todas las estaciones. Se especula que esto podría deberse a
una celda de Hadley simulada más intensificada y con menor variación
estacional. Se destaca además la ausencia de un máximo asociado al jet
subpolar, lo cual podría deberse a la falta de suficientes niveles
verticales en la porción superior.

Dada la importancia del jet polar en la dinámica atmosférica durante los
meses de invierno y primavera, su mala representación es una limitación
muy importante del modelo SPEEDY.

\subsection{Gradiente meridional de vorticidad
absoluta}\label{gradiente-meridional-de-vorticidad-absoluta}

\begin{landscape}\begin{figure}

{\centering \includegraphics[width=247mm,]{fig/tesis/etady-sp-nc-1} 

}

\caption{Gradiente meridional de vorticidad absoluta ($1\times10^{11}(ms)^{-1}$) (SPEEDY).}\label{fig:etady-sp-nc}
\end{figure}
\end{landscape}

La \autoref{fig:etady-sp-nc} muestra el campo de \(\eta_{y}\) simulado
por SPEEDY en 200hPa. Comparando con la figura \autoref{fig:etady-ncep},
se ve que la franja de máximo gradiente presente en todas las estaciones
es más zonal en el caso de SPEEDY y corrida hacia el sur en verano. La
región de gradientes negativos que se desarrolla en invierno sobre Nueva
Zelanda tiene menor extensión y no aparece en otoño ni primavera.

En la \autoref{fig:ks-sp-nc-corte} se muestra un corte meridional del
número de onda estacionario en 180°. En verano no hay regiones
prohibidas en ningún modelo, pero la región entre 45°S y 60°S aparece
como un mínimo en SPEEDY y un máximo en NCEP. Esto implica que la
propagación meridional está inhibida para un amplio rango de números de
onda en SPEEDY pero no en NCEP. En particular, la QS3 puede propagarse
meridionalmente en las observaciones, pero no en el modelo. En otoño, la
situación es inversa: alrededor de 40°S, NCEP muestra una región con
número de onda estacionario imaginario, impidiendo la propagación
meridional, mientras que en SPEEDY no existe tan impedimento. La QS3
nuevamente se ve afectada, teniendo propagación meridional irrestricta
al norte de 60°S en SPEEDY pero quedando atrapada al sur de 50°S en las
observaciones. En invierno y primavera la concordancia entre SPEEDY y
NCEP es mayor.

\begin{figure*}
\includegraphics[width=155mm, ]{fig/tesis/ks-sp-nc-corte-1} \caption{Número de onda estacionario en 200hPa en 180°O.}\label{fig:ks-sp-nc-corte}
\end{figure*}

\subsection{Función corriente}\label{funcion-corriente-2}

\begin{landscape}\begin{figure}

{\centering \includegraphics[width=247mm,]{fig/tesis/psi-sp-1} 

}

\caption{$\psi$ media en 200hPa (contornos cada $2\times10^{-7}m^2/s$), anomalía zonal de $\psi$ (sombreado,  $0,25\times10^{-7}m^2/s$) y flujos de actividad de onda medios.}\label{fig:psi-sp}
\end{figure}
\end{landscape}

La función corriente en 200hPa de SPEEDY se muestra en la
\autoref{fig:psi-sp} (donde el sombreado corresponde a la anomalía zonal
y las flechas a los flujos de actividad de onda) en comparación con los
mismos campos de NCEP (\autoref{fig:psi-ncep}, notar que en NCEP está en
coordenadas \(\sigma\)). Existe una correspondencia general buena en los
trópicos y la localización de los máximos y mínimos coincide
aproximadamente en todas las estaciones. En verano y primavera no se
simula el anticiclón de Bolivia relacionado con el SAMS, pero sí la baja
del nordeste. La intensidad de las anomalías es menor en todas las
estaciones, especialmente en el HN. Consecuentemente, también tienen
menor magnitud los flujos de actividad de onda, aunque sin embargo las
simulaciones representan dispersión media de energía de los trópicos a
los extratrópicos.

\subsection{Onda 3}\label{onda-3-1}

\begin{figure*}
\includegraphics[width=155mm, ]{fig/tesis/ampl-sp-nc-1} \caption{Amplitud de la QS3 a partir de Fourier (mgp) (SPEEDY sombreado, NCEP contornos).}\label{fig:ampl-sp-nc}
\end{figure*}

En la \autoref{fig:ampl-sp-nc} se muestra la amplitud de la QS3 media
para SPEEDY en contornos y NCEP en líneas (ambos calculados según el
método AM). En verano, entre 45°S y 60°S SPEEDY coincide con NCEP en la
localización y extensión del máximo, pero subestima su intensidad. El
máximo secundario en latitudes bajas aparece más al sur y en un nivel
más alto en SPEEDY. En otoño, el modelo carece casi por completo de
señal de QS3 en comparación con NCEP. Sin embargo, la amplitud de la QS3
simulada por SPEEDY utilizando el método MA sí es importante (no se
muestra). Esto indica que la diferencia con NCEP se debe no a la falta
de una onda 3, sino a su característica no estacionaria. En invierno, la
amplitud de la QS3 media está subestimada en SPEEDY y corrida hacia el
polo. Además, está mucho más restringida a niveles troposféricos en
comparación con NCEP; posiblemente como consecuencia de la falta de
niveles verticales altos y consecuentemente a la mala representación del
jet subpolar. En primavera, por el contrario, la señal en SPEEDY es
considerablemente mayor que en NCEP.

La distribución horizontal de Z* explicada por la QS3 simulada se
muestra en la \autoref{fig:qs3-sp-nc}. En verano, invierno y primavera,
la inclinación de los centros es hacia el este, contraria a las
observaciones e indicando transporte de cantidad de movimiento zonal
hacia el norte en vez de hacia el sur. En otoño el campo de QS3 de
SPEEDY es virtualmente nulo. En invierno y primavera, la fase de la onda
simulada por SPEEDY está defasada 180° con respecto a la observada en
NCEP.

\begin{figure*}
\includegraphics[width=155mm, ]{fig/tesis/qs3-sp-nc-1} \caption{Z* reconstruida a partir de la QS3 (mgp) (SPEEDY sombreado, NCEP contornos).}\label{fig:qs3-sp-nc}
\end{figure*}

\section{Experimentos de
sensibilidad}\label{experimentos-de-sensibilidad}

La evaluación de la corrida Control muestra que el modelo es capaz de
representar una estructura de QS3 y da evidencias del papel de los
procesos de superficie y troposféricos en mantenerla. En consecuencia,
se diseñaron y realizaron 2 corridas adicionales para explorar la
sensibilidad de la QS3 a distintos parámetros del modelo, que tienen las
siguientes características:

\begin{description}

\item[NOLAND]
 es equivalente a la corrida Control, a excepción de que se desactivaron los modelos de mar, hielo y tierra y se reemplazó en la condición de contorno, las anomalías mensuales observadas de la SST por sus correspondientes medias climatológicas. Esto implica que la simulación no considera interacción con la superficie y en cambio es solo forzada por la climatología.
 
\item[SSTZONAL]

 es equivalente a NOLAND, pero en las condiciones de contorno se consideraron las medias zonales de los promedios climatológicos mensuales de las SST. Esto implica no considerar las asimetrías zonales del forzante oceánico.

\end{description}

\subsection{Altura geopotencial}\label{altura-geopotencial-2}

Los campos de Z* en 200hPa para cada corrida se muestra en la
\autoref{fig:ghz-sp-runs}. La corrida NOLAND no presenta diferencias
importantes en la ubicación de las anomalías con respecto a Control. Se
observan algunas diferencias en la intensidad, como un debilitamiento
importante de la anomalía positiva al sudeste de Nueva Zelanda en
primavera y en la ubicación e intensidad de tal centro durante el
invierno. SSTZONAL en cambio, muestra una clara disminución de la
intensidad de las anomalías, confirmando entonces el papel de las
asimetrías de las SST medias climatológicas en determinarlas.

\begin{landscape}\begin{figure}

{\centering \includegraphics[width=247mm,]{fig/tesis/ghz-sp-runs-1} 

}

\caption{Z* en 200hPa para cada corrida de SPEEDY (mgp).}\label{fig:ghz-sp-runs}
\end{figure}
\end{landscape}

La figura \autoref{fig:ghz-dif-sp-runs} muestra la diferencia entre Z*
de la corrida Control y las corridas de sensibilidad. NOLAND sólo
muestra valores significativos en invierno y primavera. Se destaca que
la diferencia tiene una estructura de onda 3 en latitudes medias y
altas. Esto confirma que la falta de interacción entre el suelo y la
atmósfera tiene influencia sobre este patrón de onda.

La diferencia entre SSTZONAL y Control se da principalmente entre 45°S y
60°S y muestra también estructura de onda 3, en particular en primavera
y algo menos en el invierno. En cambio, en verano y otoño, el campo de
diferencia presenta evidencia de estructura de onda~1. Estos resultados
son consistentes con los de \citet{Quintanar1995} dado que la
eliminación de las asimetrías zonales de SST en los trópicos reduciría
la convección anómala que genera los trenes de onda que ellos
concluyeron son el principal sostén de este patrón.

\begin{landscape}\begin{figure}

{\centering \includegraphics[width=247mm,]{fig/tesis/ghz-dif-sp-runs-1} 

}

\caption{Diferencia de Z* (mgp) en 200hPa entre las simulaciones NOLAND y SSTZONAL y la corrida Control.}\label{fig:ghz-dif-sp-runs}
\end{figure}
\end{landscape}

\subsection{Viento zonal}\label{viento-zonal-2}

\begin{landscape}\begin{figure}

{\centering \includegraphics[width=247mm,]{fig/tesis/u-dif-sp-runs-1} 

}

\caption{Diferencia de viento zonal en 200hPa (m/s) entre las simulaciones NOLAND y SSTZONAL y la corrida Control.}\label{fig:u-dif-sp-runs}
\end{figure}
\end{landscape}

La diferencia del viento medio zonal entre corridas se muestra en la
\autoref{fig:u-dif-sp-runs}. La resta NOLAND - Control muestra valores
casi nulos en otoño. En verano NOLAND tiene vientos más intensos al
norte y sur de Nueva Zelanda y menos intensos sobre ésta región,
indicando una intensificación de los bloqueos con respeto a la corrida
Control; el mismo patrón se observa en primavera. En invierno y
primavera hay valores positivos en latitudes bajas indicando un
debilitamiento de los alisios.

En SSTZONAL el patrón de aumento de los bloqueos se ve intensificado. En
otoño, invierno y primavera se observan franjas de valores negativos al
sur de 30° y positivos al norte, indicando que el jet se encuentra
desplazado hacia el norte.

\subsection{Función corriente}\label{funcion-corriente-3}

\begin{landscape}\begin{figure}

{\centering \includegraphics[width=247mm,]{fig/tesis/psi-sp-runs-1} 

}

\caption{Anomalía zonal de $\psi$ (sombreado,  $1\times10^{-7}m^2/s$) y flujos de actividad de onda medios para cada corrida.}\label{fig:psi-sp-runs}
\end{figure}
\end{landscape}

En la \autoref{fig:psi-sp-runs} se muestra la anomalía zonal de la
función corriente para cada corrida y los flujos de acción de onda. La
misma es consistente con las observaciones generales de la
\autoref{fig:ghz-sp-runs}. En verano, las diferencias entre los flujos
de Control y NOLAND no parecen ser importantes, mientras que los flujos
han disminuido considerablemente en SSTZONAL. Esto se destaca también,
aunque en menor medida en las otras tres estaciones y confirmaría el
papel de las asimetrías zonales del forzante oceánico en mantener las
ondas cuasiestacionarias.

\subsection{Onda 3}\label{onda-3-2}

\begin{landscape}\begin{figure}

{\centering \includegraphics[width=247mm,]{fig/tesis/ampl-sp-runs-1} 

}

\caption{Amplitud de la QS3 (mgp) para cada corrida según Fourier y método AM.}\label{fig:ampl-sp-runs}
\end{figure}
\end{landscape}

En la \autoref{fig:ampl-sp-runs} se muestra la amplitud de la QS3
(método AM) para cada estación y cada corrida. Primavera es la estación
con el cambio más dramático. Mientras que la señal es muy alta en la
troposfera y latitudes medias para la corrida Control, es prácticamente
despreciable en NOLAND y SSTZONAL. En verano se observa también una
reducción de la señal de la QS3 en 60°S entre Control, NOLAND y SSTZONAL
además de un ligero corrimiento hacia el sur. La señal de más al norte,
en cambio, no presenta cambios importantes. En otoño la poca señal
existente en latitudes medias en Control y NOLAND prácticamente
desaparece en SSTZONAl pero se observa un aumento de la amplitud de la
QS3 en la costa antártica. No obstante, la baja resolución del modelo
hace cuestionable los resultados en la vecindad de una topografía tan
compleja como la Antártida.

En invierno, NOLAND tiene la mayor señal de QS3 entre las simulaciones y
Control, la mínima, con SSTZONAL presentando valores intermedios. Hay
que tener en cuenta que la falta de representación de la estratosfera
limita la simulación de las ondas estacionarias y en particular la QS3
en esta estación. En consecuencia, se requieren más estudios para
evaluar y entender este resultado que superan los planeados en esta
tesis.

Es importante notar que estas observaciones son sensibles a la
metodología utilizada. La \autoref{fig:ampl-mean-sp-runs} muestra la
amplitud media de la QS3 a partir del método MA para cada corrida y es
útil compararla con la \autoref{fig:ampl-sp-runs} (construida a partir
del método AM). Verano se comporta de manera similar, con una ligera
disminución de la señal en NOLAND y SSTZONAL pero Otoño tiene un gran
cambio. Tanto la corrida Control como NOLAND y SSTZONAL tienen una señal
fuerte que no cambia significativamente entre corridas. Esto indica que
la baja señal en otoño en SPEEDY se asocia con que la onda 3 está
presente mes a mes, pero no estacionariamente. Lo mismo parece suceder
invierno de la corrida Control y primavera de NOLAND y SSTZONAL.

\begin{landscape}\begin{figure}

{\centering \includegraphics[width=247mm,]{fig/tesis/ampl-mean-sp-runs-1} 

}

\caption{Amplitud de la QS3 (mgp) para cada corrida según Fourier y método MA.}\label{fig:ampl-mean-sp-runs}
\end{figure}
\end{landscape}

\begin{figure*}
\includegraphics[width=155mm, ]{fig/tesis/sd-fase-sp-runs-1} \caption{Desvío estándar (en grados) de la fase media mensual para cada estación y cada corrida.}\label{fig:sd-fase-sp-runs}
\end{figure*}

Para profundizar en esta observación, la \autoref{fig:sd-fase-sp-runs}
muestra el desvío estándar de la fase media mensual para cada estación y
corrida. Si se asume distribución normal, aproximadamente el 95\% de los
datos están en un rango de \(\pm 2\sigma\) al rededor de la media. Como
la fase está acotada entre 0° y 120°, valores de \(\sigma\) por encima
de 30° implica que los datos están distribuidos casi uniformemente en
todo el dominio.

En verano, el desvío es mínimo y aumenta ligeramente en NOLAND y
SSTZONAL pero siempre por debajo de los 30°. En otoño, todas las
corridas muestran valores altos muy por encima de los 30°, indicando que
hay muy poca estacionariedad en todas las corridas. En invierno, la
corrida Control está en el límite de los 30° y NOLAND y SSTZONAL tienen
valores menores. En primavera, la corrida Control tiene valores mínimos
de \(\sigma\) mientras que NOLAND y SSTZONAL presentan valores muy
altos, por encima de los 30°. Estos valores son consistentes con lo
observado en la comparación de las Figuras~\ref{fig:ampl-sp-runs}
y~\ref{fig:ampl-mean-sp-runs} y sugieren que las diferencias observadas
entre las corridas en la señal de la QS3 se deben principalmente a una
mayor o menor estacionariedad y no a una mayor o menor amplitud de las
onda planetaria 3.

\chapter{Conclusiones}\label{conclusiones}

Este trabajo describe las características de la onda planetaria de
número de onda zonal 3 (QS3) en el HS y realiza una primera exploración
sobre las posibles causas de su establecimiento y variabilidad.

Consistente con estudios previos, se encontró que la circulación media
del HS está dominada por un patrón de QS1 y la QS3 contribuye en
latitudes medias, con amplitud moderada a lo largo de las cuatro
estaciones.

El análisis de la fase de la QS3 mostró un ciclo estacional de 35° de
amplitud consistente con estudios previos pero además se observó una
considerable variabilidad interanual para cada mes. En particular, se
observó que a pesar de que Sudamérica climatológicamente se ve afectada
por un centro de anomalía anticiclónica asociado a la QS3, la
variabilidad en la fase es lo suficientemente grande como para que
exista un número no despreciable de años en los que se observa un centro
ciclónico sobre el continente. Estos resultados resaltan la importancia
de esta variable al momento de evaluar las influencias locales de la QS3
y deben tenerse en cuenta en los promedios climatológicos ya que pueden
enmascarar una señal fuerte de onda 3 mediante interferencia
destructiva.

Una exploración preliminar de la aplicación de wavelets al estudio de
las ondas planetarias resultó prometedora. Al permitir la variabilidad
de la amplitud en la dirección zonal, posibilitaría capturar variaciones
espaciales de las ondas en vez de tratarlas como características
globales de que afectan homogéneamente a todo un círculo de latitud. A
partir de esta metodología se estudió preliminarmente la estacionariedad
de las ondas planetarias y se encontró que la QS3 es más estacionaria en
la región del Índico que en el pacífico. Su estudio más profundo será
abordado en trabajos posteriores.

A través de la aplicación del análisis de componentes principales se
encontró que la estacionalidad de la QS3 no coincide exactamente con las
estaciones climatológicas clásicas. Se concluyó que la división EFM, AM,
JJ, ASO y ND representa una mejor clasificación. La misma resaltó que
durante noviembre y diciembre la QS3 no presenta rasgos de onda
estacionaria y permitió describir con mayor claridad la inclinación de
la QS3 en la vertical durante junio y julio. Esto sugiere que la
extensión vertical de la QS3 hasta la estratósfera se debe, al menos en
parte, a la propagación vertical de ondas de Rossby desde la tropósfera.

A partir de las regresiones entre las anomalías de circulación y un
índice de actividad de la QS3 definido en este trabajo (\(A_3\)) se
observó que la QS3 tiene comportamientos muy heterogéneos según el mes.
Existen meses donde la QS3 se asocia a importantes trenes de onda desde
los trópicos y otros con una estructura más zonal. Se encontró
asociación con el SAM en julio y diciembre, aunque con signo opuesto.
Esto podría indicar que la QS3 es en realidad una familia de fenómenos
que tienen en común que en una descomposición de Fourier resultan en una
onda 3 intensa.

La validación de la corrida Control realizada con el modelo SPEEDY
mostró que es capaz de generar la QS3, aunque su desempeño varía
considerablemente según la estación, con el peor desempeño en otoño. La
estructura horizontal de la QS3 simulada se caracterizó por inclinación
de los centros hacia el este en todas las estaciones, contrario a la
inclinación hacia el oeste presente en NCEP en verano, otoño y primavera
y un defasaje de 180° con respecto a la QS3 observada en NCEP.

Se observó que a pesar de no desarrollar un jet subpolar --producto de
su burda representación de la estratósfera-- el modelo reprodujo la QS1
en la tropósfera y su variación estacional, aunque crudamente. A pesar
de no carecer de inclinación con la altura --sugiriendo ausencia de
propagación vertical--, la QS1 simulada por SPEEDY alcanza máximos en
los niveles más altos disponibles.

Los experimentos de sensibilidad mostraron variaciones importantes en la
QS3, principalmente en invierno y primavera. En primavera, la amplitud
de la QS3 disminuye dramáticamente ausencia de procesos de interacción
suelo-atmosfera. Esto ocurre también en el verano, aunque en menor
medida. En invierno en cambio, eliminar la interacción entre la
superficie y la atmósfera incrementó la intensidad de la QS3 en
comparación con la corrida Control, posiblemente por la aparición de un
tren de ondas de número de onda zonal 3 con propagación meridional que
está presente en las observaciones, pero ausente en la corrida Control.
En esa estación, la eliminación de las anomalías zonales de SST en la
corrida SSTZONAL no cambió la amplitud de la QS3 significativamente en
comparación con la NOLAND. Dada la ausencia de representación de la
estratosfera en el modelo, los resultados del invierno pueden no ser
concluyentes y requieren mas estudios futuros.

La comparación entre el método AM y el MA para la caracterización de la
onda 3 mostró que el cambio en la amplitud de la QS3 entre corridas no
se debió a una disminución en la actividad de la onda planetaria 3 sino
a una reducción en su estacionariedad. Esto sugiere que la onda 3 puede
generarse a partir de la variabilidad interna de la atmósfera pero que
las influencias superficiales cumplen un rol importante en su
localización. Alternativamente, si se considera que distintos mecanismos
posiblemente independientes pueden generar una señal de onda 3, es
posible que los influenciados con las condiciones de superficie tengan
características más estacionarias que los relacionados con la
variabilidad interna de la atmósfera.

Un análisis más completo de las corridas de sensibilidad realizadas será
tema de futuras investigaciones, así como el uso de modelos más
sofisticados.

\renewcommand\bibname{Referencias}
\bibliography{Papers/Biblio.bib,Papers/packages.bib}


\end{document}
